\documentclass[a4paper,14pt]{extarticle}

\usepackage{fontspec}
\usepackage{polyglossia}
\setmainlanguage{russian}
\setotherlanguage{english}
\setmainfont[Ligatures=TeX]{Times New Roman}
\newfontfamily\cyrillicfont[Ligatures=TeX]{Times New Roman}
\newfontfamily\cyrillicfonttt{Menlo}[Scale=0.88]
\newfontfamily\cyrillicfontsf{Arial}

\usepackage[
  left=30mm,
  right=15mm,
  top=20mm,
  bottom=20mm
]{geometry}
\usepackage{setspace}
\onehalfspacing
\setlength{\parindent}{1.25cm}
\setlength{\parskip}{0pt}

% Сдаваемая версия ВКР не должна разрывать слова переносами вида
% "обозначе- / ния", особенно на стыке страниц.
\hyphenpenalty=10000
\exhyphenpenalty=10000
\brokenpenalty=10000
\doublehyphendemerits=1000000
\finalhyphendemerits=1000000

\usepackage{indentfirst}
\usepackage{microtype}
\usepackage{graphicx}
\usepackage{adjustbox}
\usepackage{float}
\usepackage{flafter}
\usepackage{caption}
\usepackage{longtable}
\usepackage{booktabs}
\usepackage{array}
\usepackage{calc}
\usepackage{ragged2e}
\usepackage{etoolbox}
\usepackage{enumitem}
\usepackage{amsmath}
\usepackage{amssymb}
\usepackage{fvextra}
\usepackage{listings}
\usepackage{lastpage}
\usepackage{titlesec}
\usepackage{tocloft}
\usepackage{needspace}
\usepackage{placeins}
\usepackage{xurl}
\usepackage{hyperref}
\usepackage{bookmark}
\usepackage{xcolor}

\graphicspath{{./}{../}{docx_pipeline/figures/png/}{docx_pipeline/figures/svg/}}

\hypersetup{
  unicode=true,
  colorlinks=false,
  pdfborder={0 0 0},
  pdftitle={Разработка и внедрение интеллектуальной системы управления микроклиматом в теплице на основе машинного обучения},
  pdfauthor={Ларионов Сергей Иванович}
}

\setcounter{secnumdepth}{0}
\setcounter{tocdepth}{3}

\titleformat{\section}
  {\normalfont\bfseries}
  {}
  {0pt}
  {}
\titlespacing*{\section}{1.25cm}{0pt}{12pt}

\titleformat{\subsection}
  {\normalfont\bfseries}
  {}
  {0pt}
  {}
\titlespacing*{\subsection}{1.25cm}{12pt}{6pt}

\titleformat{\subsubsection}
  {\normalfont\bfseries}
  {}
  {0pt}
  {}
\titlespacing*{\subsubsection}{1.25cm}{10pt}{6pt}

\pretocmd{\section}{\FloatBarrier\clearpage}{}{}
\pretocmd{\subsection}{\Needspace{6\baselineskip}}{}{}
\pretocmd{\subsubsection}{\Needspace{5\baselineskip}}{}{}

\newcommand{\VkrStructuralSection}[1]{%
  \FloatBarrier
  \clearpage
  \phantomsection
  \addcontentsline{toc}{section}{#1}%
  \begin{center}\bfseries #1\end{center}%
  \par\nopagebreak\vspace{12pt}%
}

\newcommand{\VkrStructuralSectionNoToc}[1]{%
  \begin{center}\bfseries #1\end{center}%
  \par\nopagebreak\vspace{12pt}%
}

\newcommand{\VkrAppendixSection}[2]{%
  \FloatBarrier
  \clearpage
  \phantomsection
  \addcontentsline{toc}{section}{ПРИЛОЖЕНИЕ #1 #2}%
  \begin{center}\bfseries ПРИЛОЖЕНИЕ #1\end{center}%
  \par\nopagebreak\vspace{6pt}%
  \begin{center}\bfseries #2\end{center}%
  \par\nopagebreak\vspace{12pt}%
}

\renewcommand{\cftsecleader}{\cftdotfill{\cftdotsep}}
\makeatletter
\renewcommand{\tableofcontents}{%
  \section*{СОДЕРЖАНИЕ}%
  \@starttoc{toc}%
}
\makeatother
\setlength{\cftbeforesecskip}{2pt}
\setlength{\cftbeforesubsecskip}{1pt}

\captionsetup{
  font={normalsize,stretch=1},
  labelformat=empty,
  labelsep=none,
  justification=centering,
  singlelinecheck=false,
  skip=6pt
}
\captionsetup[lstlisting]{
  font={normalsize,stretch=1},
  labelformat=simple,
  labelsep=endash,
  justification=raggedright,
  singlelinecheck=false,
  skip=6pt
}

\setlist[itemize]{label=\textemdash,leftmargin=1.25cm,itemsep=0pt,topsep=0pt,parsep=0pt}
\setlist[enumerate]{leftmargin=1.25cm,itemsep=0pt,topsep=0pt,parsep=0pt}

\lstset{
  basicstyle=\ttfamily\small,
  breaklines=true,
  columns=fullflexible,
  keepspaces=true,
  frame=single,
  framerule=0.4pt,
  framesep=4pt,
  captionpos=t,
  rulecolor=\color{black},
  xleftmargin=0pt,
  xrightmargin=0pt,
  aboveskip=6pt,
  belowskip=6pt
}
\renewcommand{\lstlistingname}{Листинг}
\RecustomVerbatimEnvironment{verbatim}{Verbatim}{
  breaklines=true,
  breakanywhere=true,
  fontsize=\small,
  frame=single,
  framerule=0.4pt,
  framesep=4pt
}
\renewcommand{\texttt}[1]{\textnormal{#1}}

\sloppy
\emergencystretch=2em
\clubpenalty=10000
\widowpenalty=10000
\displaywidowpenalty=10000
\pagestyle{plain}

\providecommand{\tightlist}{%
  \setlength{\itemsep}{0pt}\setlength{\parskip}{0pt}}
\newcommand{\VkrFigureFull}[1]{%
  \adjustbox{max width=\linewidth,max totalheight=0.74\textheight,keepaspectratio,center}{#1}%
}
\newcommand{\VkrFigureInline}[1]{%
  \adjustbox{max width=0.94\linewidth,max totalheight=0.55\textheight,keepaspectratio,center}{#1}%
}
\newcommand{\VkrFigureArchitecture}[1]{%
  \adjustbox{max width=0.98\linewidth,max totalheight=0.74\textheight,keepaspectratio,center}{#1}%
}
\newcommand{\VkrFigureTurnCycle}[1]{%
  \adjustbox{max width=0.98\linewidth,max totalheight=0.76\textheight,keepaspectratio,center}{#1}%
}
\newcommand{\VkrFigureStagePolicy}[1]{%
  \adjustbox{max width=0.98\linewidth,max totalheight=0.65\textheight,keepaspectratio,center}{#1}%
}
\newcommand{\VkrFigureExperimentDesign}[1]{%
  \adjustbox{max width=0.98\linewidth,max totalheight=0.60\textheight,keepaspectratio,center}{#1}%
}
\newcommand{\VkrFigureTall}[1]{%
  \adjustbox{max width=0.88\linewidth,max totalheight=0.50\textheight,keepaspectratio,center}{#1}%
}
\newcommand{\VkrFigureVertical}[1]{%
  \adjustbox{max width=0.88\linewidth,max totalheight=0.56\textheight,keepaspectratio,center}{#1}%
}
\newcommand{\VkrFigureCompact}[1]{%
  \adjustbox{max width=0.92\linewidth,max totalheight=0.48\textheight,keepaspectratio,center}{#1}%
}
\newcommand{\VkrFigureWideCompact}[1]{%
  \adjustbox{max width=0.88\linewidth,max totalheight=0.38\textheight,keepaspectratio,center}{#1}%
}
\providecommand{\pandocbounded}[1]{\VkrFigureFull{#1}}
\providecommand{\passthrough}[1]{#1}
\newcounter{none}

\setlength{\LTpre}{0pt}
\setlength{\LTpost}{6pt}
\setlength{\tabcolsep}{3.5pt}
\setlength{\extrarowheight}{2pt}
\renewcommand{\arraystretch}{1.22}
\AtBeginEnvironment{longtable}{\small\singlespacing}

\newcommand{\VkrTableCaption}[1]{%
  \Needspace{0.44\textheight}%
  \par\smallskip
  {\singlespacing\noindent #1\par}%
  \vspace{2pt}%
}


\begin{document}
\hypersetup{pageanchor=false}
\newcommand{\VkrTitleHeaderLine}[1]{%
  \noindent{\fontsize{12pt}{14pt}\selectfont\makebox[\textwidth][c]{#1}\par}%
}
\newcommand{\VkrUdkField}{УДК \underline{\hspace{1.8cm}}}

\begin{titlepage}
\fontsize{12pt}{14pt}\selectfont
\singlespacing
\VkrTitleHeaderLine{МИНИСТЕРСТВО СЕЛЬСКОГО ХОЗЯЙСТВА РОССИЙСКОЙ ФЕДЕРАЦИИ}

\vspace{0.55cm}

\VkrTitleHeaderLine{ФЕДЕРАЛЬНОЕ ГОСУДАРСТВЕННОЕ БЮДЖЕТНОЕ}
\VkrTitleHeaderLine{ОБРАЗОВАТЕЛЬНОЕ УЧРЕЖДЕНИЕ ВЫСШЕГО ОБРАЗОВАНИЯ}
\VkrTitleHeaderLine{«САНКТ-ПЕТЕРБУРГСКИЙ ГОСУДАРСТВЕННЫЙ}
\VkrTitleHeaderLine{АГРАРНЫЙ УНИВЕРСИТЕТ»}

\vspace{0.75cm}

\VkrTitleHeaderLine{ИНСТИТУТ ЭНЕРГЕТИКИ}

\vspace{0.75cm}

\VkrTitleHeaderLine{КАФЕДРА ТЕПЛОТЕХНИКИ И ЭНЕРГООБЕСПЕЧЕНИЯ}

\vspace{0.9cm}

\begin{flushright}
На правах рукописи

\vspace{0.2cm}

\VkrUdkField
\end{flushright}

\vspace{0.75cm}

\begin{center}
Ларионов Сергей Иванович

\vspace{0.8cm}

\textbf{«Разработка и внедрение интеллектуальной системы управления микроклиматом в теплице на основе машинного обучения для повышения эффективности выращивания растений»}

\vspace{0.85cm}

Выпускная квалификационная работа магистра

\vspace{0.65cm}

Направление подготовки\\
«Теплотехника и теплоэнергетика в агропромышленном комплексе»
\end{center}

\vspace{0.4cm}

\begin{flushright}
\begin{minipage}{0.46\textwidth}
Выпускная квалификационная\\
работа защищена

\vspace{0.15cm}

«\underline{\hspace{0.8cm}}»\underline{\hspace{2.5cm}}2026 г.

\vspace{0.18cm}

Оценка

\vspace{0.12cm}

\underline{\hspace{3.0cm}}

\vspace{0.18cm}

Секретарь ГЭК \underline{\hspace{2.2cm}}
\end{minipage}
\end{flushright}

\vfill

\begin{center}
г. Санкт-Петербург\\
2026
\end{center}
\end{titlepage}

\setcounter{page}{2}
\hypersetup{pageanchor=true}

\VkrStructuralSectionNoToc{РЕФЕРАТ}\label{ux440ux435ux444ux435ux440ux430ux442}

Работа содержит \pageref{LastPage} с., 0 рис., 8 табл., 32 источн., 1
прил.

Выпускная квалификационная работа посвящена разработке и внедрению
интеллектуальной системы управления микроклиматом в теплице на основе
машинного обучения. В работе теплица рассматривается как
теплотехническая и киберфизическая система, объединяющая сенсорные узлы,
исполнительные устройства, MQTT-брокер, веб-панель оператора, правила
автоматизации, профили управления и локальное хранилище истории.

Цель работы состоит в создании расширяемой программно-аппаратной
IoT/ML-платформы для мониторинга и управления тепловлажностным режимом
теплицы с возможностью подключения разнородных датчиков, регистрации
исполнительных воздействий, накопления данных и дальнейшего применения
методов машинного обучения для прогнозирования температуры, влажности и
концентрации CO\(_2\).

В аналитической части рассмотрены современные подходы к smart
greenhouse, IoT-архитектурам, прогнозированию микроклимата,
MPC/RL-управлению и энергоэффективности. В проектной части описана
архитектура системы на базе микроконтроллеров ESP8266, Arduino Mega
2560, MQTT-брокера Mosquitto, Flask-dashboard и SQLite. В практической
части описаны прошивки сенсорных узлов, контроллер реле, веб-панель,
журналирование данных, экспорт истории, правила и профили управления.

Результатом работы является прототип расширяемой IoT/ML-платформы,
которая обеспечивает сбор телеметрии, управление 15 релейными каналами,
хранение истории измерений и событий, а также подготовку данных для
ML-эксперимента. Дополнительно подготовлены техническая спецификация,
протокол испытаний, схемы архитектуры, приложения, список литературы и
сборка черновика в LaTeX/PDF, DOCX и HTML.

Ключевые слова: теплица, микроклимат, интеллектуальное управление,
машинное обучение, MQTT, IoT, ESP8266, Arduino Mega, Flask, SQLite,
прогнозирование временных рядов.

\subsection{Abstract}\label{abstract}

The thesis is devoted to the development and implementation of an
intelligent greenhouse microclimate control system based on machine
learning. The greenhouse is considered as both a thermal engineering
object and a cyber-physical system that combines sensor nodes,
actuators, an MQTT broker, an operator dashboard, automation rules,
control profiles, and local historical data storage.

The goal of the work is to create an extensible IoT/ML hardware and
software platform for greenhouse microclimate monitoring and control,
with support for heterogeneous sensors, actuator event logging,
historical data collection, and further machine learning-based
forecasting of temperature, humidity, and CO\(_2\) concentration.

The analytical part reviews modern approaches to smart greenhouses, IoT
architectures, microclimate forecasting, MPC/RL control, and energy
efficiency. The design part describes the system architecture based on
ESP8266 microcontrollers, Arduino Mega 2560, Mosquitto MQTT broker,
Flask dashboard, and SQLite. The implementation part describes sensor
firmware, relay controller firmware, web dashboard, data logging,
history export, automation rules, and control profiles.

The result is a prototype of an extensible IoT/ML platform that supports
telemetry collection, control of 15 relay channels, storage of sensor
and actuator history, and data preparation for an ML experiment. The
project also contains a technical specification, test protocol,
architecture diagrams, appendices, bibliography, and LaTeX/PDF, DOCX,
and HTML exports.

Keywords: greenhouse, microclimate, intelligent control, machine
learning, MQTT, IoT, ESP8266, Arduino Mega, Flask, SQLite, time series
forecasting.

\clearpage

\VkrStructuralSectionNoToc{СОДЕРЖАНИЕ}
\makeatletter\@starttoc{toc}\makeatother
\clearpage

\VkrStructuralSection{ПЕРЕЧЕНЬ СОКРАЩЕНИЙ И ОБОЗНАЧЕНИЙ}\label{ux441ux43fux438ux441ux43eux43a-ux441ux43eux43aux440ux430ux449ux435ux43dux438ux439-ux438-ux43eux431ux43eux437ux43dux430ux447ux435ux43dux438ux439}

{\def\LTcaptype{none} % do not increment counter
\begin{longtable}[]{@{}
  >{\RaggedRight\arraybackslash}p{(\linewidth - 4\tabcolsep) * \real{0.3333}}
  >{\RaggedRight\arraybackslash}p{(\linewidth - 4\tabcolsep) * \real{0.3333}}
  >{\RaggedRight\arraybackslash}p{(\linewidth - 4\tabcolsep) * \real{0.3333}}@{}}
\hline
\begin{minipage}[b]{\linewidth}\RaggedRight
Сокращение
\end{minipage} & \begin{minipage}[b]{\linewidth}\RaggedRight
Расшифровка
\end{minipage} & \begin{minipage}[b]{\linewidth}\RaggedRight
Пояснение в контексте работы
\end{minipage} \\
\hline
\endhead
\hline
\endlastfoot
АСУ & автоматизированная система управления & система мониторинга и
управления микроклиматом теплицы \\
ВКР & выпускная квалификационная работа & итоговая работа по теме
проекта \\
БД & база данных & SQLite-хранилище истории измерений и событий \\
ИИ & искусственный интеллект & интеллектуальный слой анализа и поддержки
решений \\
ТЭН & трубчатый электронагреватель & исполнительное устройство
нагрева \\
CO\(_2\) & carbon dioxide & концентрация углекислого газа в воздухе
теплицы \\
CSV & comma-separated values & текстовый формат выгрузки табличных
данных \\
DHT11 & Digital Humidity and Temperature sensor & датчик температуры и
влажности \\
ESP8266 & микроконтроллер с Wi-Fi & платформа сенсорных узлов \\
GPIO & general-purpose input/output & цифровой вход/выход
микроконтроллера \\
GRU & gated recurrent unit & рекуррентная нейросетевая архитектура \\
HTML & hypertext markup language & формат веб-интерфейса и экспортного
просмотра \\
HTTP & hypertext transfer protocol & протокол веб-доступа к dashboard \\
IoT & internet of things & интернет вещей, объединение устройств через
сеть \\
JSON & JavaScript Object Notation & формат MQTT-сообщений датчиков и
сводок \\
LSTM & long short-term memory & рекуррентная нейросетевая архитектура
для временных рядов \\
MAE & mean absolute error & метрика средней абсолютной ошибки
прогноза \\
ML & machine learning & машинное обучение \\
MPC & model predictive control & модельно-предиктивное управление \\
MQTT & message queuing telemetry transport & легкий
publish/subscribe-протокол обмена сообщениями \\
PNG & portable network graphics & формат изображения для схем и
скриншотов \\
REST & representational state transfer & стиль HTTP-интерфейсов \\
RL & reinforcement learning & обучение с подкреплением \\
RMSE & root mean squared error & корень из средней квадратичной
ошибки \\
RNN & recurrent neural network & рекуррентная нейронная сеть \\
SCD30 & CO\(_2\), humidity, and temperature sensor & датчик CO\(_2\),
температуры и влажности \\
SQL & structured query language & язык запросов к базе данных \\
SQLite & embedded SQL database & локальная база истории dashboard \\
SVG & scalable vector graphics & векторный формат схем \\
TCN & temporal convolutional network & сверточная архитектура для
временных рядов \\
UI & user interface & пользовательский интерфейс \\
XLSX & Office Open XML spreadsheet & формат Excel-выгрузки \\
\end{longtable}
}

\VkrStructuralSection{ВВЕДЕНИЕ}\label{ux432ux432ux435ux434ux435ux43dux438ux435}

Современные тепличные комплексы требуют точного контроля параметров
микроклимата: температуры, влажности, концентрации CO\(_2\), режимов
вентиляции, нагрева, увлажнения и полива. С теплотехнической точки
зрения теплица является объектом с тепловым балансом, влажностным
режимом и вентиляционным воздухообменом. Даже в малой теплице параметры
изменяются не изолированно: включение вентилятора влияет на температуру
и влажность, открытие заслонок меняет воздухообмен, работа нагревателей
имеет инерционный эффект, а концентрация CO\(_2\) зависит от вентиляции,
фотосинтетической активности растений и притока наружного воздуха.
Поэтому управление микроклиматом нельзя сводить только к ручному
включению оборудования или простым пороговым действиям без накопления
истории.

Актуальность работы связана с переходом тепличных систем от локальной
автоматики к киберфизическим IoT-решениям. В современных исследованиях
smart greenhouse рассматривается как система, объединяющая датчики,
исполнительные механизмы, сетевой обмен, хранилище данных, веб-интерфейс
и интеллектуальный аналитический слой {[}1; 2{]}. Такая постановка
особенно важна для малых и учебно-экспериментальных теплиц: если система
проектируется только под один датчик и одно реле, она быстро перестает
быть исследовательской платформой. Для дальнейшего применения машинного
обучения требуется расширяемая архитектура, в которой новые сенсорные
узлы, исполнительные каналы и расчетные признаки подключаются без
переписывания всей системы.

Отдельное значение имеет энергоэффективность. Нагрев, вентиляция,
увлажнение и освещение потребляют значительные ресурсы, а вентиляционные
воздействия могут приводить к теплопотерям {[}28; 31{]}. При этом
реакция теплицы на управляющее воздействие зависит от текущего состояния
среды, внешних условий, инерции конструкций, плотности растений и режима
воздухообмена. Поэтому управление должно учитывать не только текущие
значения датчиков, но и прогноз развития микроклимата {[}14; 25{]}.

Методы машинного обучения позволяют использовать накопленную телеметрию
для прогнозирования температуры, влажности, CO\(_2\), роста растений и
урожайности. В литературе применяются LSTM, RNN, TCN, Transformer и
вероятностные модели, ориентированные на временные ряды тепличных данных
{[}9; 10; 11; 13{]}. Однако для практического применения таких моделей
сначала необходим надежный инженерный слой: датчики должны стабильно
публиковать данные, исполнительные устройства должны управляться через
понятный интерфейс, а история измерений и воздействий должна сохраняться
в пригодном для анализа виде.

В рамках данной работы рассматривается разработка и опытное внедрение
расширяемой программно-аппаратной IoT/ML-платформы мониторинга и
управления микроклиматом теплицы. Платформа объединяет микроконтроллеры,
MQTT-брокер, веб-панель оператора, правила автоматизации, профили
управления, журналирование данных и подготовку истории для машинного
обучения. Реальная история измерений и управляющих воздействий
необходима для перехода от эмпирического управления к имитационной или
ML-модели тепловлажностного режима {[}32; 29{]}. Машинное обучение в
работе рассматривается как верхний интеллектуальный слой, который
опирается на данные, собранные инженерной системой. Такой подход
позволяет сохранить надежность базового управления и одновременно
подготовить основу для прогнозных моделей и дальнейшего
интеллектуального управления.

Цель работы: разработать и внедрить расширяемую IoT/ML-платформу
мониторинга и управления тепловлажностным режимом теплицы с возможностью
подключения разнородных датчиков, управления исполнительными
устройствами, накопления истории и последующего применения методов
машинного обучения для прогнозирования и интеллектуальной поддержки
управления.

Для достижения цели необходимо решить следующие задачи:

\begin{enumerate}
\def\labelenumi{\arabic{enumi})}
\tightlist
\item
  Проанализировать современные подходы к автоматизации микроклимата
  теплиц, IoT-архитектурам, прогнозированию микроклимата и
  интеллектуальному управлению.
\item
  Определить требования к расширяемой системе мониторинга и управления:
  состав параметров, датчики, исполнительные механизмы, канал связи,
  хранение данных, интерфейс оператора и пригодность истории для
  ML-анализа.
\item
  Спроектировать архитектуру системы на базе микроконтроллеров, MQTT,
  серверного приложения, базы данных истории и единой модели обмена
  сообщениями.
\item
  Реализовать программно-аппаратный контур: сенсорные узлы, релейный
  контроллер, MQTT-обмен, dashboard, журналирование показаний и
  состояний.
\item
  Подготовить основу для интеллектуального слоя: накопление временных
  рядов, экспорт данных, правила и профили управления, план
  ML-эксперимента.
\item
  Проверить работоспособность системы на опытном объекте, сформировать
  датасет эксплуатации и определить направления дальнейшего развития.
\end{enumerate}

Объект исследования: тепловлажностный режим теплицы как управляемая
киберфизическая система агропромышленного комплекса.

Предмет исследования: методы и программно-аппаратные средства
мониторинга, управления и интеллектуальной поддержки регулирования
микроклимата теплицы.

Практическая значимость работы состоит в создании системы, которая может
использоваться как основа для эксплуатации малой теплицы и для
дальнейших экспериментов с прогнозированием микроклимата. Реализованная
архитектура позволяет собирать телеметрию, управлять исполнительными
устройствами, хранить историю и формировать датасет для машинного
обучения.

Научная новизна работы заключается в архитектурно-методическом подходе к
построению малой теплицы как расширяемой IoT/ML-платформы. В отличие от
одноразовой схемы «датчик - реле - панель управления», в работе
рассматривается единый контур, где разнородные датчики публикуют данные
в общей модели обмена, управляющие воздействия регистрируются совместно
с параметрами микроклимата, а накопленная история становится основой для
последующего ML-прогнозирования реакции теплицы. При этом не заявляется
разработка нового алгоритма машинного обучения: интеллектуальный слой
рассматривается как прогнозная надстройка над безопасным правиловым
контуром.

Структура работы включает введение, три главы, заключение, список
литературы и приложения. В первой главе приведен аналитический обзор
источников по smart greenhouse, IoT, прогнозированию и методам
управления. Во второй главе описано проектирование системы. В третьей
главе рассмотрены реализация, проверка работоспособности и подготовка
ML-эксперимента.

\section{1 Аналитический обзор систем интеллектуального управления
микроклиматом
теплицы}\label{ux430ux43dux430ux43bux438ux442ux438ux447ux435ux441ux43aux438ux439-ux43eux431ux437ux43eux440-ux441ux438ux441ux442ux435ux43c-ux438ux43dux442ux435ux43bux43bux435ux43aux442ux443ux430ux43bux44cux43dux43eux433ux43e-ux443ux43fux440ux430ux432ux43bux435ux43dux438ux44f-ux43cux438ux43aux440ux43eux43aux43bux438ux43cux430ux442ux43eux43c-ux442ux435ux43fux43bux438ux446ux44b}

\subsection{1.1 Теплица как киберфизическая система
управления}\label{ux442ux435ux43fux43bux438ux446ux430-ux43aux430ux43a-ux43aux438ux431ux435ux440ux444ux438ux437ux438ux447ux435ux441ux43aux430ux44f-ux441ux438ux441ux442ux435ux43cux430-ux443ux43fux440ux430ux432ux43bux435ux43dux438ux44f}

Современная теплица является не просто защищенным сооружением для
выращивания растений, а киберфизической системой, в которой датчики,
исполнительные механизмы, микроконтроллеры, серверное программное
обеспечение и алгоритмы принятия решений работают как единый контур.
Целевыми параметрами такого контура выступают температура воздуха,
относительная влажность, концентрация CO\(_2\), освещенность, влажность
почвы или субстрата, режим полива и состояние растений. Поддержание этих
параметров влияет на урожайность, скорость роста, устойчивость растений
к стрессу и энергоемкость выращивания.

В обзорной литературе последних лет тепличное хозяйство все чаще
рассматривается через призму smart greenhouse, то есть управляемой
среды, где IoT, искусственный интеллект, спектральные датчики,
мультимодальное слияние данных и интеллектуальные системы используются
для повышения эффективности выращивания. В обзоре El Ouaham et
al.~показано, что развитие smart greenhouse связано не только с заменой
ручного контроля на датчики, но и с переходом к системам, которые
собирают разнородные данные, анализируют их и поддерживают принятие
решений {[}1{]}. Для данной ВКР это принципиально: ценность системы
определяется не только тем, что она измеряет температуру и включает
реле, а тем, что она создает расширяемый контур данных для последующего
прогнозирования и анализа реакции теплицы.

Русскоязычные источники подтверждают ту же постановку на инженерном
уровне. В сравнительном анализе современных систем автоматизированного
управления микроклиматом подчеркивается, что автоматизация теплиц
направлена на снижение эксплуатационных затрат, повышение урожайности и
переход от простого поддержания параметров к интеллектуальному
управлению {[}3{]}. Халина и Цыбанюк описывают базовую структуру
автоматизированной системы: датчики измеряют параметры среды, блок
управления сравнивает их с заданными значениями, после чего
исполнительные механизмы изменяют температуру, влажность, освещение или
другие параметры {[}4{]}. Такая схема является исходной, но в
современных условиях она требует расширения прогнозными и аналитическими
функциями.

Особенность тепличного микроклимата состоит в инерционности и
взаимосвязанности параметров. Нагрев, вентиляция, увлажнение, подача
CO\(_2\) и освещение влияют друг на друга, а реакция растений
проявляется не мгновенно. Тарков и Сукьясов рассматривают задачу
управления параметрами микроклимата через уравнения баланса и отмечают
проблему запаздывания процесса регулирования {[}5{]}. Это означает, что
реактивное управление по текущему отклонению от уставки не всегда
обеспечивает оптимальный результат: система может опоздать с
воздействием или создать избыточные энергетические затраты.

\subsection{1.2 Тепловлажностный режим и энергетическая рамка
теплицы}\label{ux442ux435ux43fux43bux43eux432ux43bux430ux436ux43dux43eux441ux442ux43dux44bux439-ux440ux435ux436ux438ux43c-ux438-ux44dux43dux435ux440ux433ux435ux442ux438ux447ux435ux441ux43aux430ux44f-ux440ux430ux43cux43aux430-ux442ux435ux43fux43bux438ux446ux44b}

Для направления «Теплотехника и теплоэнергетика в агропромышленном
комплексе» принципиально важно рассматривать теплицу как объект тепло- и
массообмена. Температура, влажность и CO\(_2\) формируются под действием
солнечной радиации, теплопотерь через ограждения, вентиляционного
воздухообмена, испарения и транспирации растений, нагрева, увлажнения и
работы вспомогательного оборудования. В обзоре Chen et al.~тепличная
система прямо описывается как сложная многовходовая и многовыходовая
система, где управление температурой, влажностью, светом, CO\(_2\) и
воздушными потоками должно одновременно обеспечивать рост растений и
энергосбережение {[}28{]}.

Динамические модели микроклимата теплицы обычно опираются на баланс
энергии, энтальпии и массы. Это важно для текущей ВКР: даже если
практическая реализация использует датчики и реле, физическая сущность
задачи остается теплотехнической. Самарин и др. формулируют задачу
имитационного моделирования оптимального управления микроклиматом
сельскохозяйственного объекта через температурно-влажностный режим,
тепловой баланс, влажностное состояние конструкций и энергосберегающий
режим выращивания {[}32{]}. Такой подход поддерживает выбранную рамку:
система должна не только включать устройства, но и формировать данные
для модели поведения теплицы.

Отдельная проблема состоит в пространственной неоднородности
микроклимата. Ogunlowo et al.~показывают, что распределение тепла и
массы в одно- и многопролетных теплицах влияет на точность оценки
энергетической потребности; при игнорировании распределения параметров
возможна переоценка или недооценка расчетов {[}29{]}. Это обосновывает
применение нескольких измерительных узлов внутри теплицы и последующий
анализ расхождений между датчиками, а не только использование одного
среднего значения.

Вентиляция является одним из ключевых каналов воздействия на
микроклимат. Она одновременно снижает температуру, удаляет влагу,
изменяет воздухообмен и влияет на концентрацию CO\(_2\). Li et al.~на
данных девяти односкатных теплиц показывают, что открытие вентиляции
меняет температуру, относительную влажность, VPD и CO\(_2\), а также
зависит от плотности посадки, высоты растений, внутренних покрытий и
времени открытия вентиляционного проема {[}30{]}. Поэтому события
включения вентиляторов и приточных каналов в собственной системе должны
журналироваться вместе с температурой, влажностью и CO\(_2\): без этого
невозможно корректно анализировать реакцию теплицы на управляющее
воздействие.

Энергетический аспект особенно заметен при управлении вентиляцией и
отоплением. Ghaderi et al.~отмечают, что вентиляционные потери тепла
являются одним из значимых факторов энергетической эффективности теплиц,
а климат-контроль может составлять основную долю энергопотребления
тепличного хозяйства {[}31{]}. Следовательно, интеллектуальный слой
управления должен стремиться не только удерживать параметры в допустимом
диапазоне, но и снижать лишние переключения, преждевременные теплопотери
и компенсирующие действия после запаздывания.

\subsection{1.3 IoT-архитектура и расширяемость тепличных
систем}\label{iot-ux430ux440ux445ux438ux442ux435ux43aux442ux443ux440ux430-ux438-ux440ux430ux441ux448ux438ux440ux44fux435ux43cux43eux441ux442ux44c-ux442ux435ux43fux43bux438ux447ux43dux44bux445-ux441ux438ux441ux442ux435ux43c}

Практическое внедрение интеллектуального управления начинается с
надежного сбора телеметрии. IoT-подходы позволяют связать датчики,
исполнительные устройства, контроллеры и программные сервисы в единую
инфраструктуру. Bersani et al.~рассматривают smart greenhouse в
контексте Industry 4.0 и показывают, что IoT-системы для теплиц включают
сенсорные сети, каналы связи, удаленный мониторинг и управляющие контуры
{[}2{]}. Такой подход хорошо согласуется с текущим проектом, где уже
присутствуют прошивки микроконтроллеров и серверный dashboard.

В инженерных реализациях часто используются микроконтроллеры, Raspberry
Pi, Arduino/ESP, MQTT, веб-интерфейсы и базы данных. Кельсин в ВКР по
автоматизированной системе управления тепличным хозяйством рассматривает
архитектуру на микроконтроллерах, близкую к практическим учебным и малым
производственным системам {[}7{]}. Статья по MQTT-based smart greenhouse
monitoring дополнительно показывает, почему MQTT удобен для телеметрии:
это легкий протокол обмена сообщениями, пригодный для IoT-устройств с
ограниченными ресурсами {[}23{]}. Для текущей ВКР это важно, поскольку
коммуникационный слой должен быть достаточно простым для
микроконтроллеров и достаточно надежным для накопления данных.

Для перехода от локальной автоматики к платформе недостаточно просто
добавить веб-интерфейс. Система должна иметь расширяемую модель
подключаемых устройств: разные датчики могут измерять температуру,
влажность, CO\(_2\), освещенность, влажность субстрата или
дополнительные расчетные параметры, но для верхнего уровня они должны
попадать в историю в сопоставимой структуре. Такой подход согласуется с
направлением multimodal data fusion в smart greenhouse: данные от разных
источников становятся не разрозненными показаниями, а общей
информационной базой для анализа состояния растений и микроклимата
{[}1{]}. В практической архитектуре это означает необходимость
идентификатора устройства, типа метрики, времени измерения, единиц
измерения и признака качества данных.

Расширяемость важна и для исполнительных устройств. В малой теплице реле
может управлять вентилятором, заслонкой, насосом, увлажнителем или
нагревателем, но для анализа микроклимата все эти действия должны
фиксироваться как события воздействия на тепловлажностный режим. Без
совместной регистрации датчиков и реле невозможно отличить естественное
изменение температуры от реакции на вентиляцию или нагрев. Поэтому
журнал исполнительных воздействий является не вспомогательной функцией
dashboard, а обязательной частью ML-ready датасета.

Прикладные smart farming эксперименты демонстрируют, что тепличная
система должна рассматриваться как полный стек: датчики, исполнительные
механизмы, логика управления, хранение данных и интерфейс оператора.
Joni et al.~описывают IoT-enhanced greenhouse design для выращивания
риса, где IoT-инфраструктура связывает параметры среды и агротехническую
задачу {[}24{]}. Несмотря на отличие культуры, работа полезна как пример
построения комплексной системы. Thwin et al.~идут еще дальше и связывают
прогноз микроклимата с web integration и adaptive equipment control
{[}16{]}. Это особенно важно для данной ВКР, так как собственный проект
уже содержит dashboard, прошивки и историю событий, а значит может быть
описан как первый инкремент интеллектуального контура, а не как
завершенная одноцелевая автоматика.

Отдельно следует отметить работы, где IoT-инфраструктура связывается с
reinforcement learning. Platero-Horcajadas et al.~рассматривают
интеграцию IoT и RL для оптимизированного climate control {[}21{]}. Для
текущей ВКР этот источник важен не как требование немедленно реализовать
RL-агента, а как подтверждение архитектурного тезиса: интеллектуальный
регулятор возможен только там, где есть непрерывная телеметрия,
наблюдаемые исполнительные воздействия и воспроизводимый обмен данными.

\subsection{1.4 Прогнозирование микроклимата, роста и
урожайности}\label{ux43fux440ux43eux433ux43dux43eux437ux438ux440ux43eux432ux430ux43dux438ux435-ux43cux438ux43aux440ux43eux43aux43bux438ux43cux430ux442ux430-ux440ux43eux441ux442ux430-ux438-ux443ux440ux43eux436ux430ux439ux43dux43eux441ux442ux438}

Накопление телеметрии само по себе не делает систему интеллектуальной.
Интеллектуальный слой появляется тогда, когда исторические данные
используются для прогнозирования будущего состояния микроклимата, роста
растений или урожайности. Alhnaity et al.~применяют LSTM для
прогнозирования роста и урожайности в greenhouse environments и
показывают, что последовательностные модели способны извлекать полезные
зависимости из временных рядов тепличных данных {[}9{]}. Gong et
al.~предлагают связку TCN и RNN для прогнозирования урожайности и
сравнивают ее с классическими ML-подходами {[}10{]}. Эти работы
обосновывают использование истории наблюдений, а не только текущего
значения датчика.

Для задачи управления микроклиматом особенно важны работы, где
прогнозируются непосредственно температура, влажность и CO\(_2\). Ahn et
al.~сравнивают transformer-based и RNN-based модели для предсказания
температуры, относительной влажности и концентрации CO\(_2\) в теплице
{[}11{]}. Huang et al.~используют Attention CNN-LSTM для прогнозирования
среды грибной теплицы, также работая с температурой, влажностью и
CO\(_2\) {[}12{]}. Эти источники показывают, что микроклимат является
естественным объектом для time-series prediction, а не только для
порогового контроля. При этом качество прогноза зависит не только от
выбора модели, но и от состава признаков: полезными могут быть лаги
датчиков, время суток, сезонность, состояние вентиляции, нагрева и
увлажнения, а также различия между точками измерения.

Отдельное направление связано с прогнозным управлением и
энергоэффективностью. Aborujilah et al.~рассматривают forecast-driven
climate control для smart greenhouse и используют LSTM-модель для
proactive adjustments, направленных на снижение энергетических затрат
{[}14{]}. Soumo et al.~описывают data-driven greenhouse climate
regulation при выращивании салата с использованием BiLSTM и GRU
predictive control {[}15{]}. Thwin et al.~предлагают multivariate
multistep LSTM, который прогнозирует микроклимат и взаимодействует с
системой управления оборудованием через веб-интеграцию {[}16{]}. В
совокупности эти работы позволяют обосновать, что прогнозный слой может
быть включен в практическую АСУ теплицы.

При этом прогнозы не следует рассматривать как абсолютно надежные. Choi
and Yang предлагают probabilistic deep learning framework для greenhouse
microclimate prediction с учетом time-varying uncertainty и covariance
analysis {[}13{]}. Для диплома этот источник важен методологически: даже
если в прототипе применяется более простая модель, в тексте необходимо
признавать неопределенность прогноза и необходимость защитных
ограничений нижнего уровня управления.

\subsection{1.5 Методы управления: PID, fuzzy, MPC, RL и эвристическая
оптимизация}\label{ux43cux435ux442ux43eux434ux44b-ux443ux43fux440ux430ux432ux43bux435ux43dux438ux44f-pid-fuzzy-mpc-rl-ux438-ux44dux432ux440ux438ux441ux442ux438ux447ux435ux441ux43aux430ux44f-ux43eux43fux442ux438ux43cux438ux437ux430ux446ux438ux44f}

Классические системы управления теплицей часто строятся на пороговых
правилах, PID-регуляторах или нечеткой логике. Их преимущество состоит в
простоте, интерпретируемости и возможности реализации на
микроконтроллерах. Воробьева и др. рассматривают интеллектуальную
систему контроля температуры в теплице и сравнивают PID-регулирование с
нечетким регулятором {[}6{]}. Для инженерного проекта такие методы
остаются важными, потому что нижний уровень управления должен быть
предсказуемым и безопасным.

Однако при усложнении задачи возникает потребность в методах, которые
учитывают будущую динамику, ограничения и стоимость ресурсов. Model
Predictive Control является одним из основных подходов в научной
литературе по greenhouse climate control, поскольку он использует модель
системы для оптимизации управляющих воздействий на горизонте прогноза.
Mallick et al.~показывают, что MPC явно учитывает физические
ограничения, но сильно зависит от точности модели; для компенсации
неопределенности они предлагают связку MPC и reinforcement learning
{[}17{]}. Msaad et al.~также рассматривают RL-guided MPC для автономного
управления теплицей {[}18{]}.

Сравнение MPC и RL важно для корректного позиционирования ВКР. Morcego
et al.~рассматривают MPC как model-based подход, а RL как learning-based
подход к greenhouse climate control {[}19{]}. MPC лучше интерпретируется
и естественно работает с ограничениями, тогда как RL способен улучшать
стратегию на основе данных, но требует осторожности из-за рисков
обучения и безопасности. Поэтому для практической ВКР разумно описывать
RL/MPC не как обязательную реализацию в полном объеме, а как верхний
интеллектуальный слой, который может выдавать рекомендации или
корректировать уставки, тогда как нижний контур остается защищенным.

Перспективным направлением является включение оператора или агронома в
цикл обучения. Xiao et al.~предлагают grower-in-the-loop interactive
reinforcement learning для greenhouse climate control {[}20{]}. Такой
подход хорошо соответствует практической логике внедрения: система может
не заменять оператора полностью, а поддерживать его решения и
использовать экспертную обратную связь.

Помимо MPC/RL, в литературе встречаются эвристические оптимизационные
методы. Jawad et al.~используют artificial bee colony algorithm для
energy optimization and plant comfort management в smart greenhouse
{[}22{]}. Этот источник показывает, что задача управления микроклиматом
может решаться не только нейросетями или MPC, но и другими
оптимизационными методами. Для текущей работы это полезно как обзор
альтернатив, но основная логика ВКР остается связанной с IoT,
прогнозированием и гибридным управлением.

\subsection{1.6 Энергоэффективность и устойчивость тепличного
управления}\label{ux44dux43dux435ux440ux433ux43eux44dux444ux444ux435ux43aux442ux438ux432ux43dux43eux441ux442ux44c-ux438-ux443ux441ux442ux43eux439ux447ux438ux432ux43eux441ux442ux44c-ux442ux435ux43fux43bux438ux447ux43dux43eux433ux43e-ux443ux43fux440ux430ux432ux43bux435ux43dux438ux44f}

Энергопотребление является одним из ключевых ограничений тепличного
производства. Отопление, вентиляция, увлажнение, освещение и подача
CO\(_2\) требуют ресурсов, а реактивное управление может приводить к
избыточным включениям исполнительных механизмов. Aborujilah et al.~прямо
указывают на проблему реактивных стратегий и предлагают forecast-driven
подход с LSTM для снижения энергетических потерь {[}14{]}. Эту же линию
поддерживает теплотехническая литература по вентиляционным потерям и
интеграции отопления, охлаждения и вентиляции {[}31{]}.

Более широкий энергетический контекст раскрывается в обзоре Wu and Fu по
Agricultural Energy Internet {[}25{]}. В нем greenhouse environmental
control рассматривается как энергетически управляемый процесс наряду с
ирригацией, дополнительным освещением, логистикой и
сельскохозяйственными машинами. Для ВКР этот источник полезен в разделе
актуальности: интеллектуальное управление микроклиматом не только
улучшает условия выращивания, но и связано с задачей рационального
использования энергии.

Энергоэффективность также связана с архитектурой управления. Если
система умеет прогнозировать состояние микроклимата, она может заранее
выбирать более мягкие воздействия, снижать частоту резких переключений и
избегать компенсирующих действий после запаздывания. Однако такие
преимущества возможны только при наличии надежной телеметрии,
устойчивого обмена данными и защитных правил нижнего уровня.

\subsection{1.7 Выводы по аналитическому
обзору}\label{ux432ux44bux432ux43eux434ux44b-ux43fux43e-ux430ux43dux430ux43bux438ux442ux438ux447ux435ux441ux43aux43eux43cux443-ux43eux431ux437ux43eux440ux443}

Проведенный обзор показывает, что тема интеллектуального управления
микроклиматом теплицы находится на пересечении нескольких направлений:
теплотехники, тепло- и массообмена, IoT-архитектур, автоматизированного
управления, машинного обучения, прогнозирования временных рядов,
энергоэффективности и человеко-ориентированного принятия решений.
Практические инженерные источники описывают датчики, контроллеры, MQTT,
dashboard и исполнительные механизмы, но часто ограничиваются пороговыми
или простыми регуляторами. Современные научные работы предлагают LSTM,
RNN, Transformer, Attention CNN-LSTM, вероятностные модели, MPC и RL,
однако не всегда доводят эти методы до конкретной малой теплицы с
реальной прошивкой, пользовательским интерфейсом и совместным
журналированием измерений и воздействий.

На этом фоне текущая ВКР может быть обоснована как разработка и опытное
внедрение расширяемой IoT/ML-платформы управления микроклиматом, которая
создает основу для интеллектуального слоя. Нижний уровень системы должен
обеспечивать сбор данных, безопасное управление и связь с
исполнительными устройствами. Верхний уровень должен накапливать
телеметрию, анализировать временные ряды, формировать прогнозы и
поддерживать корректировку уставок. Научный разрыв, закрываемый работой,
состоит не в создании нового ML-алгоритма, а в переходе от разрозненной
автоматики к воспроизводимой платформе: разнородные датчики подключаются
через единую модель обмена, события реле регистрируются вместе с
микроклиматом, а история становится пригодной для ML-прогнозирования
реакции теплицы. Такой гибридный подход согласует надежность
практической АСУ с современными направлениями ML/MPC/RL и позволяет
постепенно развивать систему без риска нереалистичного обещания
полностью автономного AI-управления.

\section{2 Проектирование программно-аппаратной системы мониторинга и
управления микроклиматом
теплицы}\label{ux43fux440ux43eux435ux43aux442ux438ux440ux43eux432ux430ux43dux438ux435-ux43fux440ux43eux433ux440ux430ux43cux43cux43dux43e-ux430ux43fux43fux430ux440ux430ux442ux43dux43eux439-ux441ux438ux441ux442ux435ux43cux44b-ux43cux43eux43dux438ux442ux43eux440ux438ux43dux433ux430-ux438-ux443ux43fux440ux430ux432ux43bux435ux43dux438ux44f-ux43cux438ux43aux440ux43eux43aux43bux438ux43cux430ux442ux43eux43c-ux442ux435ux43fux43bux438ux446ux44b}

\subsection{2.1 Назначение и требования к
системе}\label{ux43dux430ux437ux43dux430ux447ux435ux43dux438ux435-ux438-ux442ux440ux435ux431ux43eux432ux430ux43dux438ux44f-ux43a-ux441ux438ux441ux442ux435ux43cux435}

Разрабатываемая система предназначена для мониторинга параметров
микроклимата теплицы, дистанционного управления исполнительными
устройствами и накопления исторических данных, пригодных для
последующего анализа и построения моделей машинного обучения. В отличие
от изолированного регулятора, такая система должна обеспечивать
одновременно три функции: сбор измерений, оперативное управление
оборудованием и сохранение наблюдений в форме, пригодной для
исследований.

По результатам анализа литературы в главе 1 можно выделить несколько
требований к архитектуре. Во-первых, система должна иметь модульную
структуру, так как современные тепличные установки часто объединяют
датчики, исполнительные механизмы, локальный сервер и пользовательский
интерфейс через сетевой протокол {[}2; 23{]}. Во-вторых, канал обмена
данными должен поддерживать малые сообщения и устойчивую работу при
нестабильной связи, что соответствует типичной практике применения MQTT
в IoT-системах теплиц {[}23; 24{]}. В-третьих, данные должны сохраняться
во временной ряд, поскольку методы прогнозирования температуры,
влажности, CO\(_2\) и урожайности требуют регулярной истории наблюдений
{[}11; 14; 13{]}.

Для дипломного проекта выбрана распределенная архитектура с отдельными
узлами измерения, контроллером реле, MQTT-брокером и веб-панелью
управления. Такой вариант лучше соответствует реальному объекту, чем
монолитная схема на одном контроллере: датчики можно размещать в разных
зонах, исполнительные каналы можно обслуживать отдельной платой, а
серверная часть может одновременно выполнять функции интерфейса
оператора, журнала событий и слоя правил.

\subsection{2.2 Общая
архитектура}\label{ux43eux431ux449ux430ux44f-ux430ux440ux445ux438ux442ux435ux43aux442ux443ux440ux430}

Система состоит из четырех основных уровней:

\begin{enumerate}
\def\labelenumi{\arabic{enumi})}
\tightlist
\item
  Уровень измерения: узлы на ESP8266 с датчиками DHT11 и SCD30 публикуют
  значения температуры, влажности и концентрации CO\(_2\).
\item
  Уровень исполнительных механизмов: контроллер Arduino Mega 2560 с
  модулем ESP8266 принимает команды по MQTT и управляет 15 релейными
  каналами.
\item
  Коммуникационный уровень: MQTT-брокер Mosquitto принимает сообщения
  датчиков, состояния реле и команды управления.
\item
  Прикладной уровень: Flask-приложение отображает состояние теплицы,
  публикует команды, выполняет правила, управляет вентиляционными
  профилями и сохраняет историю в SQLite.
\end{enumerate}

Практическая реализация этой структуры находится в каталоге
\texttt{Greenhouse/}. Серверная часть описана в
\texttt{Greenhouse/dashboard/app.py}, контейнерная инфраструктура - в
\texttt{Greenhouse/docker-compose.yml}, прошивки датчиков - в
\texttt{Greenhouse/firmware/sensor\_th/sensor\_th.ino} и
\texttt{Greenhouse/firmware/sensor\_co2/sensor\_co2.ino}, прошивка
контроллера реле - в
\texttt{Greenhouse/firmware/greenhouse\_mega/greenhouse\_mega.ino}.

При финальной сборке ВКР в этот раздел целесообразно вставить рисунок
``Общая архитектура системы'' из \texttt{thesis/diagrams.md}.

В качестве транспорта между узлами выбран MQTT. Его применение оправдано
тем, что все участники системы обмениваются короткими событиями: датчик
публикует очередное измерение, контроллер реле публикует состояние
канала, веб-сервер публикует команду включения или выключения. Такая
модель проще для расширения, чем прямые HTTP-запросы между
микроконтроллерами, поскольку новые подписчики могут получать данные без
изменения прошивок существующих узлов.

\subsection{2.3 MQTT-модель обмена
данными}\label{mqtt-ux43cux43eux434ux435ux43bux44c-ux43eux431ux43cux435ux43dux430-ux434ux430ux43dux43dux44bux43cux438}

В системе используется единое пространство топиков с префиксом
\texttt{greenhouse}. Сенсорные узлы публикуют измерения в топики вида
\texttt{greenhouse/env/\textless{}device\textgreater{}/state}, где
\texttt{\textless{}device\textgreater{}} - идентификатор устройства.
Узел температуры и влажности использует идентификатор \texttt{th-1}, а
узел CO\(_2\) - \texttt{co2-1}. Сообщения передаются в формате JSON. Для
\texttt{th-1} используются поля \texttt{ts}, \texttt{device},
\texttt{t}, \texttt{h}; для \texttt{co2-1} используются поля
\texttt{ts}, \texttt{device}, \texttt{co2}, \texttt{t}, \texttt{h}.

Состояние доступности устройств передается через LWT-топики
\texttt{greenhouse/\textless{}device\textgreater{}/status}. При
подключении устройство публикует \texttt{online}, а при аварийном
отключении брокер может опубликовать \texttt{offline}. Такой механизм
важен для практической теплицы, так как отсутствие данных может означать
не стабильное состояние среды, а отказ датчика или связи.

Управление исполнительными каналами организовано через топики
\texttt{greenhouse/relay/\textless{}n\textgreater{}/set}, где
\texttt{\textless{}n\textgreater{}} - номер реле от 1 до 15. Команды
имеют текстовый формат \texttt{ON}, \texttt{OFF} или \texttt{TOGGLE};
прошивка также допускает значения \texttt{1} и \texttt{0}. Контроллер
публикует подтвержденное состояние каждого канала в
\texttt{greenhouse/relay/\textless{}n\textgreater{}/state} со значением
\texttt{ON} или \texttt{OFF}. Кроме того, используется сводный топик
\texttt{greenhouse/relay/summary}, где передаются диагностические
сведения: время работы, уровень RSSI, состояние Wi-Fi и MQTT, а также
массив состояний реле.

Такое разделение команд и состояний снижает риск рассинхронизации
интерфейса и физического оборудования. Веб-панель не только меняет
локальное состояние кнопки, но и подписывается на подтвержденные
сообщения контроллера. Это особенно важно для тепличного оборудования,
где ошибочное отображение состояния вентилятора, насоса или нагревателя
может привести к неверным действиям оператора.

Для иллюстрации обмена сообщениями можно использовать рисунок
``MQTT-обмен сообщениями'' из \texttt{thesis/diagrams.md}.

\subsection{2.4 Узлы измерения
микроклимата}\label{ux443ux437ux43bux44b-ux438ux437ux43cux435ux440ux435ux43dux438ux44f-ux43cux438ux43aux440ux43eux43aux43bux438ux43cux430ux442ux430}

В проекте предусмотрены два типа сенсорных узлов.

Первый узел реализован в
\texttt{Greenhouse/firmware/sensor\_th/sensor\_th.ino} и использует
ESP8266 с датчиком DHT11. Он измеряет температуру и относительную
влажность и публикует данные каждые 10 секунд. Такой узел выполняет
базовую функцию контроля воздушной среды и может использоваться как
минимальный датчик для правил вентиляции или нагрева.

Второй узел реализован в
\texttt{Greenhouse/firmware/sensor\_co2/sensor\_co2.ino} и использует
ESP8266 с датчиком SCD30. Он публикует концентрацию CO\(_2\),
температуру и влажность каждые 5 секунд. В прошивке предусмотрена
фильтрация некорректных значений CO\(_2\): показания меньше 0 ppm или
больше 10000 ppm отбрасываются, после чего выполняется попытка
восстановления работы датчика. Это повышает качество исходного
временного ряда и снижает вероятность записи очевидно ошибочных
наблюдений.

Оба узла используют одинаковый принцип восстановления связи: проверяют
подключение к Wi-Fi, затем подключение к MQTT-брокеру, после чего
публикуют текущие измерения. При успешном MQTT-подключении узлы
публикуют статус \texttt{online} и, если доступно последнее измерение,
повторно отправляют его. Такая логика делает систему устойчивее к
кратковременным разрывам сети.

\subsection{2.5 Контроллер исполнительных
устройств}\label{ux43aux43eux43dux442ux440ux43eux43bux43bux435ux440-ux438ux441ux43fux43eux43bux43dux438ux442ux435ux43bux44cux43dux44bux445-ux443ux441ux442ux440ux43eux439ux441ux442ux432}

Исполнительный уровень реализован на Arduino Mega 2560 с ESP8266 в
режиме AT. Прошивка
\texttt{Greenhouse/firmware/greenhouse\_mega/greenhouse\_mega.ino}
управляет 15 релейными каналами, подключенными к цифровым выводам 22,
24, 26, 28, 30, 32, 34, 36, 38, 40, 42, 44, 46, 48 и 50. В текущей
конфигурации релейные каналы сопоставлены с насосами, вентиляторами,
заслонками, ТЭНами и УФ-фильтрами.

Контроллер подписывается на \texttt{greenhouse/relay/+/set} и
\texttt{greenhouse/relay/+/get}. Команда \texttt{set} меняет состояние
выбранного канала или всех каналов, если вместо номера используется
\texttt{all}. Команда \texttt{get} инициирует публикацию текущего
состояния без изменения реле. После подключения контроллер публикует
состояния всех реле и диагностическую сводку, что позволяет веб-панели
восстановить актуальную картину после перезапуска.

В прошивке предусмотрены периодическая публикация
\texttt{greenhouse/relay/summary}, watchdog и повторные попытки
подключения к Wi-Fi и MQTT. Эти элементы важны для автономной системы:
исполнительный контроллер должен не только выполнить команду, но и
сообщить серверу, что он остается доступным.

\subsection{2.6 Веб-панель и серверная
логика}\label{ux432ux435ux431-ux43fux430ux43dux435ux43bux44c-ux438-ux441ux435ux440ux432ux435ux440ux43dux430ux44f-ux43bux43eux433ux438ux43aux430}

Прикладной уровень реализован как Flask-приложение в
\texttt{Greenhouse/dashboard/app.py}. Оно выполняет несколько ролей:

\begin{enumerate}
\def\labelenumi{\arabic{enumi})}
\tightlist
\item
  Подписывается на MQTT-топики датчиков, реле и статусов устройств.
\item
  Хранит текущее состояние системы в памяти для быстрого отображения в
  интерфейсе.
\item
  Публикует MQTT-команды при ручном управлении реле и группами.
\item
  Выполняет пользовательские правила автоматизации.
\item
  Сохраняет историю измерений и переключений реле в SQLite.
\item
  Предоставляет веб-страницы управления, правил, профилей вентиляции и
  статистики.
\end{enumerate}

Flask-приложение получает параметры через переменные окружения.
MQTT-брокер задается переменными \texttt{MQTT\_HOST} и
\texttt{MQTT\_PORT}; файлы правил, имен реле, профилей и базы данных
задаются через \texttt{RULES\_FILE}, \texttt{NAMES\_FILE},
\texttt{PROFILES\_FILE} и \texttt{DB\_FILE}. В контейнерной конфигурации
\texttt{Greenhouse/docker-compose.yml} приложение запускается вместе с
Mosquitto и использует каталог \texttt{Greenhouse/dashboard/data} как
постоянное хранилище.

Пользовательский интерфейс содержит страницы \texttt{/},
\texttt{/rules}, \texttt{/profiles} и \texttt{/stats}. Главная страница
отображает датчики, состояние контроллера и кнопки реле. Страница правил
позволяет создавать условия включения и выключения устройств. Страница
профилей задает последовательности для вентиляции, где сначала
открываются заслонки, затем после задержки запускаются основные
устройства. Страница статистики показывает графики и позволяет
экспортировать данные в CSV или Excel.

\subsection{2.7 Правила и профили
управления}\label{ux43fux440ux430ux432ux438ux43bux430-ux438-ux43fux440ux43eux444ux438ux43bux438-ux443ux43fux440ux430ux432ux43bux435ux43dux438ux44f}

В серверной части реализованы два механизма автоматизации: правила и
профили.

Правило состоит из условия, расписания и действия. Условие задает
устройство, поле измерения (\texttt{t}, \texttt{h} или \texttt{co2}),
оператор сравнения и пороговое значение. Для предотвращения частых
переключений предусмотрен гистерезис. Расписание задает дни недели и
интервал времени; поддерживаются интервалы, переходящие через полночь.
Действие может включать или выключать отдельные реле либо запускать
профиль.

Профиль описывает последовательность команд для группы устройств. В
текущей реализации профиль содержит предварительные действия
\texttt{pre}, задержку \texttt{delay} и основные действия \texttt{main}.
Для вентиляции это позволяет сначала открыть заслонки, затем включить
вентиляторы или нагреватели. При отключении профиля порядок меняется:
сначала выключаются основные устройства, затем после задержки
выполняется обратное действие для предварительной группы. Такая логика
отражает физическую зависимость исполнительных механизмов и снижает риск
работы вентиляторов при закрытых заслонках.

С точки зрения методов управления эта часть является базовым правиловым
контроллером. В литературе подобный уровень часто рассматривается как
практический фундамент, на который затем могут накладываться MPC, RL или
прогнозные алгоритмы {[}17; 5; 4{]}. Поэтому в рамках дипломного проекта
правила и профили можно рассматривать не только как рабочую
автоматизацию, но и как интерфейс для будущего подключения моделей
машинного обучения.

Логику правил и профилей можно сопровождать рисунком ``Контур правил и
профилей'' из \texttt{thesis/diagrams.md}.

\subsection{2.8 Накопление данных для машинного
обучения}\label{ux43dux430ux43aux43eux43fux43bux435ux43dux438ux435-ux434ux430ux43dux43dux44bux445-ux434ux43bux44f-ux43cux430ux448ux438ux43dux43dux43eux433ux43e-ux43eux431ux443ux447ux435ux43dux438ux44f}

Исторические данные сохраняются в SQLite-базе \texttt{history.db}.
Таблица \texttt{sensor\_log} содержит время измерения, идентификатор
устройства, температуру, влажность и CO\(_2\). Таблица
\texttt{relay\_log} содержит время переключения, номер реле и состояние.
Для ускорения выборок по времени создаются индексы
\texttt{idx\_sensor\_ts} и \texttt{idx\_relay\_ts}.

Сенсорная история записывается периодическим фоновым потоком раз в 60
секунд. Такой шаг уменьшает объем базы по сравнению с записью каждого
MQTT-сообщения, но сохраняет регулярный временной ряд для анализа.
Переключения реле записываются при изменении состояния, что позволяет
затем сопоставлять динамику микроклимата с управляющими воздействиями.

Для последующего машинного обучения структура данных имеет несколько
полезных свойств. Во-первых, значения разделены по устройствам, поэтому
можно анализировать разные источники отдельно. Во-вторых, временная
метка есть и у измерений, и у воздействий, что позволяет формировать
признаки с лагами, скользящими средними и бинарными индикаторами
включенных исполнительных устройств. В-третьих, экспорт через
\texttt{/api/export/\textless{}data\_type\textgreater{}} позволяет
выгружать как сенсорную историю, так и журнал реле в CSV или Excel для
предварительной обработки.

Ограничение текущей структуры состоит в том, что в базе пока нет внешних
факторов: наружной температуры, солнечной радиации, прогноза погоды,
фазы роста растений и целевых агрономических показателей. В работах по
прогнозированию и оптимальному управлению такие признаки часто
существенно повышают качество моделей {[}9; 11; 14{]}. Поэтому текущий
журнал следует рассматривать как первый слой датасета, который может
быть расширен дополнительными источниками.

\subsection{2.9 Выводы по
главе}\label{ux432ux44bux432ux43eux434ux44b-ux43fux43e-ux433ux43bux430ux432ux435}

Во второй главе спроектирована и описана архитектура
программно-аппаратной системы для мониторинга и управления микроклиматом
теплицы. Система построена по распределенному принципу: датчики и
исполнительный контроллер обмениваются сообщениями через MQTT-брокер, а
веб-сервер выполняет функции интерфейса, правил автоматизации и журнала
данных.

Предложенная структура соответствует задачам дипломной работы. Она уже
обеспечивает ручное и правиловое управление, сохраняет историю измерений
и переключений, а также формирует основу для дальнейшего применения
методов машинного обучения. Следующий этап работы связан с реализацией,
испытанием и анализом накопленных данных, включая подготовку признаков
для прогнозирования микроклимата и оценки управляющих воздействий.

\section{3 Реализация и проверка программно-аппаратной
системы}\label{ux440ux435ux430ux43bux438ux437ux430ux446ux438ux44f-ux438-ux43fux440ux43eux432ux435ux440ux43aux430-ux43fux440ux43eux433ux440ux430ux43cux43cux43dux43e-ux430ux43fux43fux430ux440ux430ux442ux43dux43eux439-ux441ux438ux441ux442ux435ux43cux44b}

\subsection{3.1 Состав реализованной
системы}\label{ux441ux43eux441ux442ux430ux432-ux440ux435ux430ux43bux438ux437ux43eux432ux430ux43dux43dux43eux439-ux441ux438ux441ux442ux435ux43cux44b}

Практическая часть работы реализована как распределенная IoT-система для
теплицы. В нее входят сенсорные узлы, контроллер исполнительных
устройств, MQTT-брокер, веб-панель оператора и локальное хранилище
истории. Такой состав соответствует архитектуре, описанной во второй
главе, и позволяет рассматривать разработку не как отдельную прошивку
микроконтроллера, а как полный контур мониторинга и управления.

Реализация размещена в каталоге \texttt{Greenhouse/}. Прошивки
микроконтроллеров находятся в \texttt{Greenhouse/firmware/}, серверная
часть - в \texttt{Greenhouse/dashboard/app.py}, контейнерная
конфигурация - в \texttt{Greenhouse/docker-compose.yml}. В репозитории
также есть файл \texttt{Greenhouse/Greenhouse/Схема\ щитка.md}, который
содержит ссылку на изображение схемы щитка. Для текстового описания в
данной главе используется фактическая привязка реле из прошивки и
веб-панели, так как она является проверяемым исходным кодом проекта.

Для иллюстрации состава системы можно использовать рисунок ``Общая
архитектура системы'' из \texttt{thesis/diagrams.md}.

Основная инженерная особенность реализации состоит в разделении функций
между устройствами. Датчики не управляют реле напрямую, а только
публикуют телеметрию. Релейный контроллер не принимает решений по
микроклимату, а выполняет команды и подтверждает состояние каналов.
Серверная часть объединяет данные, интерфейс, правила и журналирование.
Такое разделение упрощает отладку и соответствует подходам IoT-систем
тепличного мониторинга {[}2; 23{]}.

\subsection{3.2 Реализация сенсорных
узлов}\label{ux440ux435ux430ux43bux438ux437ux430ux446ux438ux44f-ux441ux435ux43dux441ux43eux440ux43dux44bux445-ux443ux437ux43bux43eux432}

Узел температуры и влажности реализован в файле
\texttt{Greenhouse/firmware/sensor\_th/sensor\_th.ino}. Он построен на
ESP8266 и датчике DHT11, подключенном к GPIO2. Прошивка подключается к
Wi-Fi-сети \texttt{Greenhouse}, затем к MQTT-брокеру
\texttt{192.168.1.112:1883}, формирует топики
\texttt{greenhouse/env/th-1/state} и \texttt{greenhouse/th-1/status} и
публикует измерения каждые 10 секунд.

Сообщение сенсорного узла имеет JSON-структуру:

\begin{verbatim}
{"ts":123456,"device":"th-1","t":24.5,"h":61.0}
\end{verbatim}

Если чтение DHT11 возвращает некорректное значение, публикация
пропускается. Это предотвращает попадание пустых или нечисловых значений
в общий поток телеметрии. При восстановлении MQTT-соединения узел
публикует статус \texttt{online} и повторно отправляет последнее
валидное измерение, если оно уже было получено.

Узел CO\(_2\) реализован в файле
\texttt{Greenhouse/firmware/sensor\_co2/sensor\_co2.ino}. Он использует
ESP8266 и датчик SCD30 по интерфейсу I2C: SDA подключен к D2/GPIO4, SCL
- к D1/GPIO5. Узел публикует данные в
\texttt{greenhouse/env/co2-1/state} каждые 5 секунд и использует
LWT-топик \texttt{greenhouse/co2-1/status}.

Сообщение узла CO\(_2\) имеет вид:

\begin{verbatim}
{"ts":123456,"device":"co2-1","co2":820,"t":24.7,"h":60.5}
\end{verbatim}

В прошивке реализована проверка диапазона CO\(_2\): значения меньше 0
ppm и больше 10000 ppm отбрасываются. Если SCD30 недоступен, узел
пытается переинициализировать датчик. Такая логика важна для будущего
машинного обучения: модель прогноза чувствительна к выбросам, поэтому
первичная фильтрация на уровне прошивки уменьшает объем последующей
очистки данных.

\subsection{3.3 Реализация контроллера
реле}\label{ux440ux435ux430ux43bux438ux437ux430ux446ux438ux44f-ux43aux43eux43dux442ux440ux43eux43bux43bux435ux440ux430-ux440ux435ux43bux435}

Контроллер исполнительных устройств реализован в
\texttt{Greenhouse/firmware/greenhouse\_mega/greenhouse\_mega.ino}.
Аппаратная платформа - Arduino Mega 2560 с ESP8266 в режиме AT.
Контроллер управляет 15 релейными каналами, подключенными к выводам 22,
24, 26, 28, 30, 32, 34, 36, 38, 40, 42, 44, 46, 48 и 50.

В прошивке задано пространство MQTT-топиков с базовым префиксом
\texttt{greenhouse}. Контроллер подписывается на
\texttt{greenhouse/relay/+/set} и \texttt{greenhouse/relay/+/get}. При
получении команды \texttt{set} он меняет состояние одного реле или всех
реле, если в топике указан \texttt{all}. Поддерживаются команды
\texttt{ON}, \texttt{OFF}, \texttt{TOGGLE}, а также числовые варианты
\texttt{1} и \texttt{0}. После изменения состояния контроллер публикует
подтверждение в
\texttt{greenhouse/relay/\textless{}n\textgreater{}/state}.

Для диагностики используется топик \texttt{greenhouse/relay/summary}. В
него периодически публикуются время работы, счетчики переподключений
Wi-Fi и MQTT, RSSI и состояние всех реле. При подключении контроллер
публикует LWT-статус \texttt{greenhouse/mega-1/status} со значением
\texttt{online}, а в случае аварийного отключения брокер может отметить
устройство как \texttt{offline}.

В текущей веб-панели реле имеют прикладные имена: фильтры, насосы,
вентиляторы, заслонки и ТЭНы. Группы исполнительных устройств заданы на
сервере: заслонки - 14, 12, 11, 9; вентиляторы - 10, 8; ТЭНы - 15, 13;
насосы - 4, 6, 7, 5; фильтры - 3, 2, 1. Это позволяет оператору работать
не только с отдельными номерами реле, но и с функциональными группами
оборудования.

\subsection{3.4 Реализация веб-панели и серверной
части}\label{ux440ux435ux430ux43bux438ux437ux430ux446ux438ux44f-ux432ux435ux431-ux43fux430ux43dux435ux43bux438-ux438-ux441ux435ux440ux432ux435ux440ux43dux43eux439-ux447ux430ux441ux442ux438}

Серверная часть реализована как Flask-приложение. Оно запускается в
контейнере \texttt{dashboard}, подключается к MQTT-брокеру
\texttt{mosquitto} и использует каталог \texttt{/data} для постоянных
файлов. Основные параметры задаются переменными окружения:
\texttt{MQTT\_HOST}, \texttt{MQTT\_PORT}, \texttt{RULES\_FILE},
\texttt{NAMES\_FILE}, \texttt{DB\_FILE}, \texttt{PROFILES\_FILE},
\texttt{DASHBOARD\_PASSWORD} и \texttt{TZ\_OFFSET}.

При подключении к MQTT сервер подписывается на:

\begin{itemize}
\tightlist
\item
  \texttt{greenhouse/relay/+/state};
\item
  \texttt{greenhouse/mega-1/status};
\item
  \texttt{greenhouse/relay/summary};
\item
  \texttt{greenhouse/env/+/state};
\item
  \texttt{greenhouse/+/status}.
\end{itemize}

Полученные сообщения обновляют оперативное состояние приложения.
Сенсорные сообщения сохраняются в структуре
\texttt{state{[}"sensors"{]}}, состояния реле - в
\texttt{state{[}"relays"{]}}, диагностическая сводка Mega - в
\texttt{state{[}"mega\_summary"{]}}. Такой подход позволяет быстро
отображать последние значения в интерфейсе, не обращаясь к базе данных
при каждом обновлении страницы.

Пользовательский интерфейс содержит четыре основных раздела. Главная
страница \texttt{/} показывает состояние контроллера, датчиков, групп и
отдельных реле. Страница \texttt{/rules} используется для настройки
правил автоматизации. Страница \texttt{/profiles} задает
последовательности вентиляции и нагрева с задержками. Страница
\texttt{/stats} отображает историю и предоставляет экспорт данных.

Ручное управление реализовано через HTTP-маршруты
\texttt{/relay/\textless{}num\textgreater{}/toggle},
\texttt{/relay/all/on}, \texttt{/relay/all/off},
\texttt{/group/\textless{}name\textgreater{}/on} и
\texttt{/group/\textless{}name\textgreater{}/off}. Эти маршруты не
управляют GPIO напрямую, а публикуют MQTT-команды в
\texttt{greenhouse/relay/\textless{}n\textgreater{}/set}. Поэтому
веб-панель остается верхним уровнем управления, а физическое
переключение выполняется только контроллером реле.

\subsection{3.5 Журналирование и экспорт
данных}\label{ux436ux443ux440ux43dux430ux43bux438ux440ux43eux432ux430ux43dux438ux435-ux438-ux44dux43aux441ux43fux43eux440ux442-ux434ux430ux43dux43dux44bux445}

Для хранения истории используется SQLite. При старте приложение создает
таблицы \texttt{sensor\_log} и \texttt{relay\_log}, а также индексы по
времени. Таблица \texttt{sensor\_log} хранит время, устройство,
температуру, влажность и CO\(_2\). Таблица \texttt{relay\_log} хранит
время, номер реле и состояние.

Структуру таблиц можно сопровождать ER-рисунком ``SQLite-структура
истории'' из \texttt{thesis/diagrams.md}.

Сенсорные данные записываются фоновым потоком раз в 60 секунд. Это
значение выбрано как компромисс между частотой исходных MQTT-сообщений и
размером базы. Датчик CO\(_2\) публикует данные каждые 5 секунд, датчик
температуры и влажности - каждые 10 секунд, но для долгосрочного анализа
микроклимата минутный шаг часто достаточен. Переключения реле
записываются как события изменения состояния, что позволяет позже
сопоставлять управляющие воздействия и реакцию микроклимата.

Экспорт реализован через маршрут
\texttt{/api/export/\textless{}data\_type\textgreater{}}. Для
\texttt{sensors} выгружаются строки из \texttt{sensor\_log}, для
\texttt{relays} - строки из \texttt{relay\_log}. Поддерживаются
HTML-просмотр, CSV и Excel. Это делает систему пригодной не только для
эксплуатации, но и для подготовки экспериментального датасета.

\subsection{3.6 Реализация правил и
профилей}\label{ux440ux435ux430ux43bux438ux437ux430ux446ux438ux44f-ux43fux440ux430ux432ux438ux43b-ux438-ux43fux440ux43eux444ux438ux43bux435ux439}

Правила автоматизации сохраняются в JSON-файл \texttt{rules.json}.
Каждое правило может содержать условие по датчику, расписание, список
действий с реле или ссылку на профиль. Условие включает устройство, поле
(\texttt{t}, \texttt{h} или \texttt{co2}), оператор сравнения, порог и
гистерезис. Расписание задается днями недели и интервалом времени,
включая интервалы с переходом через полночь.

Фоновый поток правил выполняет проверку примерно каждые 5 секунд. Если
условие истинно и расписание разрешает действие, сервер публикует
команды реле или запускает профиль. Если у правила включен режим
\texttt{reverse}, при переходе условия в ложное состояние выполняется
обратное действие. Такой механизм соответствует базовому пороговому
управлению и подходит как надежный нижний слой, который в будущем можно
дополнить прогнозными рекомендациями {[}4; 17{]}.

Профили вентиляции сохраняются в \texttt{profiles.json}. Профиль состоит
из предварительных действий \texttt{pre}, задержки \texttt{delay},
основных действий \texttt{main} и задержки отключения
\texttt{delay\_off}. Типичный сценарий: сначала открыть заслонки, затем
после задержки включить вентиляторы или ТЭНы. При остановке профиля
сначала отключаются основные устройства, затем после задержки
закрываются заслонки. Такая последовательность учитывает физическую
логику оборудования и уменьшает риск некорректного порядка операций.

\subsection{3.7 Проверка
работоспособности}\label{ux43fux440ux43eux432ux435ux440ux43aux430-ux440ux430ux431ux43eux442ux43eux441ux43fux43eux441ux43eux431ux43dux43eux441ux442ux438}

Проверка работоспособности системы должна подтверждать не отдельный
запуск программы, а замкнутый контур обмена данными и команд. Для этого
используются следующие критерии:

\begin{enumerate}
\def\labelenumi{\arabic{enumi})}
\tightlist
\item
  Сенсорные узлы подключаются к Wi-Fi и MQTT-брокеру и публикуют статусы
  \texttt{online}.
\item
  Узел \texttt{th-1} публикует JSON-сообщения с температурой и
  влажностью в \texttt{greenhouse/env/th-1/state}.
\item
  Узел \texttt{co2-1} публикует JSON-сообщения с CO\(_2\), температурой
  и влажностью в \texttt{greenhouse/env/co2-1/state}.
\item
  Веб-панель получает сообщения датчиков и отображает последние
  значения.
\item
  Команда из веб-панели публикуется в
  \texttt{greenhouse/relay/\textless{}n\textgreater{}/set}.
\item
  Контроллер Mega изменяет состояние реле и публикует подтверждение в
  \texttt{greenhouse/relay/\textless{}n\textgreater{}/state}.
\item
  Веб-панель обновляет состояние реле по подтвержденному MQTT-сообщению.
\item
  Измерения записываются в \texttt{sensor\_log}, а переключения реле - в
  \texttt{relay\_log}.
\item
  Страница статистики выгружает историю в CSV или Excel.
\item
  Профиль выполняет последовательность
  \texttt{pre\ -\textgreater{}\ delay\ -\textgreater{}\ main}, а при
  остановке - обратную последовательность.
\end{enumerate}

Часть этих критериев проверяется статически по исходному коду, часть
требует запуска оборудования. В рамках репозитория подтверждаются
структура MQTT-топиков, обработчики Flask-маршрутов, создание таблиц
SQLite, логика правил, логика профилей и функции экспорта. Для финальной
защиты желательно дополнить главу скриншотами dashboard, примером
MQTT-сообщений из брокера и фрагментом выгруженного CSV-файла.

\subsection{3.8 Подготовка к
ML-эксперименту}\label{ux43fux43eux434ux433ux43eux442ux43eux432ux43aux430-ux43a-ml-ux44dux43aux441ux43fux435ux440ux438ux43cux435ux43dux442ux443}

Текущая реализация уже формирует основу для интеллектуального слоя: в
базе появляются временные ряды температуры, влажности и CO\(_2\), а
также события включения и выключения исполнительных устройств. На первом
этапе разумно решать задачу краткосрочного прогноза микроклимата, а не
задачу полностью автономного управления. Такой выбор согласуется с
литературой, где прогнозирование микроклиматических параметров часто
рассматривается как предварительный этап для adaptive control и MPC/RL
{[}11; 14; 16; 13{]}.

Для эксперимента можно сформировать датасет с минутной дискретизацией.
Целевыми переменными будут температура, влажность и CO\(_2\) через
заданный горизонт прогноза, например 10, 30 или 60 минут. В качестве
признаков используются текущие и лаговые значения датчиков, время суток,
день недели, состояние реле, агрегаты за последние интервалы и признаки
недавних переключений. Базовыми метриками качества следует выбрать MAE и
RMSE; для сравнения нужна простая baseline-модель, например прогноз
последним наблюдением или скользящим средним.

Пайплайн подготовки данных и оценки моделей можно показать рисунком
``Накопление данных и ML-пайплайн'' из \texttt{thesis/diagrams.md}.

Такой эксперимент не требует немедленного внедрения ML-модели в контур
управления. Его цель - доказать, что собранная системой телеметрия
пригодна для прогнозирования и что архитектура проекта позволяет перейти
от правил к интеллектуальной поддержке решений.

\subsection{3.9 Выводы по
главе}\label{ux432ux44bux432ux43eux434ux44b-ux43fux43e-ux433ux43bux430ux432ux435-1}

В третьей главе описана фактическая реализация системы мониторинга и
управления микроклиматом теплицы. Реализованы сенсорные узлы ESP8266,
контроллер реле на Arduino Mega 2560, MQTT-обмен, Flask-dashboard,
правила автоматизации, профили вентиляции, SQLite-журнал и экспорт
данных.

Полученная система закрывает инженерную часть дипломной работы и создает
основу для дальнейшего ML-слоя. Ее сильная сторона состоит в том, что
данные и управляющие события проходят через единый контур MQTT и
сохраняются в форме, пригодной для анализа. Основные дальнейшие работы
связаны с накоплением достаточного объема телеметрии, очисткой временных
рядов, построением baseline-прогнозов и последующим сравнением с
ML-моделями.

\VkrStructuralSection{ЗАКЛЮЧЕНИЕ}\label{ux437ux430ux43aux43bux44eux447ux435ux43dux438ux435}

В ходе работы была подготовлена и описана программно-аппаратная система
мониторинга и управления микроклиматом теплицы. Система построена как
распределенный IoT-контур: сенсорные узлы публикуют телеметрию,
контроллер реле управляет исполнительными устройствами, MQTT-брокер
обеспечивает обмен сообщениями, а серверная часть предоставляет
веб-интерфейс, правила автоматизации, профили управления и хранение
истории.

В аналитической части рассмотрены современные направления автоматизации
теплиц: smart greenhouse, IoT-архитектуры, MQTT-обмен, прогнозирование
микроклимата, методы MPC/RL и энергоэффективное управление. Обзор
показал, что практическая ценность интеллектуального управления
появляется только при наличии надежного слоя сбора данных и исполнения
команд. Поэтому в работе выбран поэтапный подход: сначала реализуется
инженерная система мониторинга и управления, затем на ее данных может
строиться ML-слой.

В проектной части определены требования к системе и описана архитектура,
включающая сенсорный слой, исполнительный контроллер, MQTT-брокер,
Flask-dashboard и SQLite-хранилище. Выбранная архитектура разделяет
функции измерения, управления, хранения и интерфейса. Это упрощает
расширение системы и позволяет добавлять новые датчики, исполнительные
каналы или аналитические модули без полной переработки существующих
прошивок.

В практической части описана реализация основных компонентов. Узел
\texttt{sensor\_th} измеряет температуру и влажность, узел
\texttt{sensor\_co2} измеряет CO\(_2\), температуру и влажность,
контроллер \texttt{greenhouse\_mega} управляет 15 релейными каналами, а
веб-панель принимает MQTT-сообщения, отображает состояние теплицы и
публикует команды управления. В серверной части реализованы правила
автоматизации, профили вентиляции с задержками, журналирование сенсорных
данных и событий реле, а также экспорт истории в CSV и Excel.

Результаты работы:

\begin{enumerate}
\def\labelenumi{\arabic{enumi})}
\tightlist
\item
  Сформирован корпус литературы по smart greenhouse, IoT, MQTT,
  ML-прогнозированию и интеллектуальному управлению микроклиматом.
\item
  Подготовлен аналитический обзор, обосновывающий связь инженерной АСУ
  теплицы с будущим ML-слоем.
\item
  Спроектирована архитектура системы мониторинга и управления
  микроклиматом.
\item
  Описана фактическая реализация сенсорных узлов, релейного контроллера,
  MQTT-обмена, dashboard и SQLite-журнала.
\item
  Зафиксирована техническая спецификация MQTT-топиков, JSON-сообщений,
  HTTP-маршрутов, релейных групп и структуры базы данных.
\item
  Подготовлен план ML-эксперимента для прогноза температуры, влажности и
  CO\(_2\) на основе \texttt{sensor\_log} и \texttt{relay\_log}.
\end{enumerate}

Поставленная цель достигнута на уровне создания инженерной основы
интеллектуальной системы: реализован контур сбора данных, управления
оборудованием, хранения истории и подготовки телеметрии для машинного
обучения. При этом работа не делает необоснованного вывода о полностью
автономном AI-управлении: текущая реализация является надежной базой,
поверх которой можно проводить прогнозные эксперименты и постепенно
вводить интеллектуальные рекомендации.

Основные ограничения текущей версии связаны с набором данных. В системе
пока отсутствуют наружная температура, солнечная радиация, освещенность,
фаза роста культуры и фактические агрономические показатели. Кроме того,
для обучения устойчивых моделей требуется накопить достаточно длинную
историю наблюдений в разных режимах эксплуатации. Эти ограничения не
снижают ценность инженерной части, но задают границы интерпретации
будущих ML-результатов.

Дальнейшее развитие работы целесообразно вести по следующим
направлениям:

\begin{enumerate}
\def\labelenumi{\arabic{enumi})}
\tightlist
\item
  Накопить датасет за длительный период эксплуатации теплицы.
\item
  Добавить внешние признаки: погоду, освещенность, наружную влажность и
  режимы выращивания.
\item
  Реализовать пайплайн подготовки данных из SQLite в обучающий датасет.
\item
  Сравнить baseline-прогнозы с табличными ML-моделями и
  последовательностными моделями.
\item
  Добавить режим рекомендаций для оператора на основе прогноза.
\item
  Рассмотреть интеграцию прогнозного слоя с правилами или MPC-подобной
  логикой управления.
\end{enumerate}

Таким образом, выполненная работа создает основу для дальнейшего
перехода от ручного и правилового управления к data-driven подходу, где
решения об управлении микроклиматом опираются на историю наблюдений,
прогноз и анализ реакции теплицы на управляющие воздействия.

\VkrStructuralSection{СПИСОК ИСПОЛЬЗОВАННЫХ ИСТОЧНИКОВ}\label{ux441ux43fux438ux441ux43eux43a-ux438ux441ux43fux43eux43bux44cux437ux43eux432ux430ux43dux43dux44bux445-ux438ux441ux442ux43eux447ux43dux438ux43aux43eux432}

\begin{enumerate}
\def\labelenumi{\arabic{enumi}.}
\tightlist
\item
  El Ouaham W., Sadik M., Ennajih A., Mouzouna Y., Orchi H., Elouaham S.
  Smart Greenhouses in the Era of IoT and AI: A Comprehensive Review of
  AI Applications, Spectral Sensing, Multimodal Data Fusion, and
  Intelligent Systems // Agriculture. 2026. Vol. 16, No.~7. Article 761.
  DOI: 10.3390/agriculture16070761.
\item
  Bersani C., Ruggiero C., Sacile R., Soussi A., Zero E. Internet of
  Things Approaches for Monitoring and Control of Smart Greenhouses in
  Industry 4.0 // Energies. 2022. Vol. 15, No.~10. Article 3834. DOI:
  10.3390/en15103834.
\item
  Ван Вэйянь, Харисов А. Р. Сравнительный анализ современных систем
  автоматизированного управления микроклиматом в теплице // Вестник
  науки. 2025. № 12 (93). С. 256-275. URL:
  \url{https://cyberleninka.ru/article/n/sravnitelnyy-analiz-sovremennyh-sistem-avtomatizirovannogo-upravleniya-mikroklimatom-v-teplitse.}
\item
  Халина Т. М., Цыбанюк Д. Е. Разработка автоматизированной системы
  управления микроклиматом тепличного комплекса // Наука и молодежь:
  материалы XVII Всероссийской научно-технической конференции. Барнаул:
  АлтГТУ, 2020. С. 125-126.
\item
  Тарков Ю. М., Сукьясов С. В. Построение модели управления параметрами
  микроклимата в теплице // Материалы Иркутского ГАУ. 2025. С. 836-839.
\item
  Воробьева Н. С. {[}и др.{]}. Интеллектуальная система контроля
  температуры в теплице // Известия Нижневолжского агроуниверситетского
  комплекса. 2026. № 2 (86). С. 517-522.
\item
  Кельсин А. А. Автоматизированная система управления тепличным
  хозяйством на базе микроконтроллеров: выпускная квалификационная
  работа. Екатеринбург: Уральский государственный педагогический
  университет, 2021.
\item
  Гречиха А. С. Обзор системы управления микроклиматом
  автоматизированной теплицы для выращивания микрозелени // Агротехника
  и энергообеспечение. 2023. № 2.
\item
  Alhnaity B., Pearson S., Leontidis G., Kollias S. Using Deep Learning
  to Predict Plant Growth and Yield in Greenhouse Environments.
  arXiv:1907.00624. 2019. URL: \url{https://arxiv.org/abs/1907.00624.}
\item
  Gong L., Yu M., Jiang S., Cutsuridis V., Pearson S. Deep Learning
  Based Prediction on Greenhouse Crop Yield Combined TCN and RNN //
  Sensors. 2021. Vol. 21, No.~13. Article 4537. DOI: 10.3390/s21134537.
\item
  Ahn J. Y., Kim Y., Park H., Park S. H., Suh H. K. Evaluating
  Time-Series Prediction of Temperature, Relative Humidity, and CO\(_2\)
  in the Greenhouse with Transformer-Based and RNN-Based Models //
  Agronomy. 2024. Vol. 14, No.~3. Article 417. DOI:
  10.3390/agronomy14030417.
\item
  Huang S., Liu Q., Wu Y., Chen M., Yin H., Zhao J. Edible Mushroom
  Greenhouse Environment Prediction Model Based on Attention CNN-LSTM //
  Agronomy. 2024. Vol. 14, No.~3. Article 473. DOI:
  10.3390/agronomy14030473.
\item
  Choi W.-J., Yang M. Probabilistic Deep Learning Framework for
  Greenhouse Microclimate Prediction with Time-Varying Uncertainty and
  Covariance Analysis // Agriculture. 2025. Vol. 15, No.~23. Article
  2461. DOI: 10.3390/agriculture15232461.
\item
  Aborujilah A., Al-Sarem M., Abu-Zanona M. A. Forecast-Driven Climate
  Control for Smart Greenhouses: Energy Optimization Using LSTM Model //
  Energies. 2025. Vol. 18, No.~21. Article 5821. DOI:
  10.3390/en18215821.
\item
  Soumo E. A., Calisti M., Polydoros A. Data-Driven Greenhouse Climate
  Regulation in Lettuce Cultivation Using BiLSTM and GRU Predictive
  Control. arXiv:2507.21669. 2026. URL:
  \url{https://arxiv.org/abs/2507.21669.}
\item
  Thwin K. M. M., Horanont T., Phatrapornnant T. Machine-Learning
  Microclimate Forecasting for Adaptive Equipment Control via Web
  Integration in Open-Ventilated Greenhouses // AgriEngineering. 2024.
  Vol. 6, No.~3. P. 2845-2869. DOI: 10.3390/agriengineering6030165.
\item
  Mallick S., Airaldi F., Dabiri A., Sun C., De Schutter B.
  Reinforcement learning-based model predictive control for greenhouse
  climate control // Smart Agricultural Technology. 2025. Vol. 10.
  Article 100751. DOI: 10.1016/j.atech.2024.100751.
\item
  Msaad S., Harraway M., McAllister R. D. RL-Guided MPC for Autonomous
  Greenhouse Control. arXiv:2506.13278. 2025. URL:
  \url{https://arxiv.org/abs/2506.13278.}
\item
  Morcego B., Yin W., Boersma S., van Henten E., Puig V., Sun C.
  Reinforcement Learning Versus Model Predictive Control on Greenhouse
  Climate Control. arXiv:2303.06110. 2023. URL:
  \url{https://arxiv.org/abs/2303.06110.}
\item
  Xiao M., Lan J., Yu J., Ma W., Xie Q., Sun C. Grower-in-the-Loop
  Interactive Reinforcement Learning for Greenhouse Climate Control.
  arXiv:2505.23355. 2025. URL: \url{https://arxiv.org/abs/2505.23355.}
\item
  Platero-Horcajadas M., Pardo-Pina S., Cámara-Zapata J.-M.,
  Brenes-Carranza J.-A., Ferrández-Pastor F.-J. Enhancing Greenhouse
  Efficiency: Integrating IoT and Reinforcement Learning for Optimized
  Climate Control // Sensors. 2024. Vol. 24, No.~24. Article 8109. DOI:
  10.3390/s24248109.
\item
  Jawad M., Wahid F., Ali S., Ma Y., Alkhyyat A., Khan J., Lee Y. Energy
  optimization and plant comfort management in smart greenhouses using
  the artificial bee colony algorithm // Scientific Reports. 2025. DOI:
  10.1038/s41598-024-84141-5.
\item
  Teja G. S., Sathish P. MQTT Protocol based Smart Greenhouse
  Environment Monitoring System using Machine Learning // International
  Journal of Innovative Technology and Exploring Engineering. 2020. Vol.
  9, No.~9. P. 278-283. DOI: 10.35940/ijitee.I7149.079920.
\item
  Joni I. M. {[}et al.{]}. Smart Farming Experiment: IoT-Enhanced
  Greenhouse Design for Rice Cultivation with Foliar and Soil
  Fertilization // AgriEngineering. 2025. Vol. 7, No.~11. Article 380.
  DOI: 10.3390/agriengineering7110380.
\item
  Wu Y., Fu X. Clean and Smart Energy Technologies for Agricultural
  Energy Internet Systems: A Comprehensive Review and Future
  Perspectives // Applied Sciences. 2026. Vol. 16, No.~10. Article 4859.
  DOI: 10.3390/app16104859.
\item
  Hindi I., Alsharkawi A., Al-Ajlouni M., Qarallah B. Enhancing
  autonomous agriculture control systems in greenhouses for sustainable
  resource usage using deep learning techniques // PLOS One. 2026. DOI:
  10.1371/journal.pone.0344946.
\item
  Тюхтий Ю. А. Разработка интеллектуальной системы управления
  температурой в помещении: магистерская диссертация. Екатеринбург:
  Уральский федеральный университет, 2021.
\item
  Chen S., Liu A., Tang F., Hou P., Lu Y., Yuan P. A Review of
  Environmental Control Strategies and Models for Modern Agricultural
  Greenhouses // Sensors. 2025. Vol. 25, No.~5. Article 1388. DOI:
  10.3390/s25051388.
\item
  Ogunlowo Q. O., Akpenpuun T. D., Na W.-H., Rabiu A., Adesanya M. A.,
  Addae K. S., Kim H.-T., Lee H.-W. Analysis of Heat and Mass
  Distribution in a Single- and Multi-Span Greenhouse Microclimate //
  Agriculture. 2021. Vol. 11, No.~9. Article 891. DOI:
  10.3390/agriculture11090891.
\item
  Li H., Li A., Hou Y., Zhang C., Guo J., Li J., Ma Y., Wang T., Yin Y.
  Analysis of Heat and Humidity in Single-Slope Greenhouses with Natural
  Ventilation // Buildings. 2023. Vol. 13, No.~3. Article 606. DOI:
  10.3390/buildings13030606.
\item
  Ghaderi M., Reddick C., Sorin M. A Systematic Heat Recovery Approach
  for Designing Integrated Heating, Cooling, and Ventilation Systems for
  Greenhouses // Energies. 2023. Vol. 16, No.~14. Article 5493. DOI:
  10.3390/en16145493.
\item
  Самарин Г. Н., Попов А. Н., Плотников С. В., Ружьев В. А. Методика
  построения имитационной модели оптимального управления параметрами
  микроклимата сельскохозяйственного объекта // Аграрный научный журнал.
  2025. № 1. С. 126-135. DOI: 10.28983/asj.y2025i1pp126-135.
\end{enumerate}

\VkrAppendixSection{А}{Фрагменты прошивки контроллера
реле}\label{ux43fux440ux438ux43bux43eux436ux435ux43dux438ux435-ux430-ux444ux440ux430ux433ux43cux435ux43dux442ux44b-ux43fux440ux43eux448ux438ux432ux43aux438-ux43aux43eux43dux442ux440ux43eux43bux43bux435ux440ux430-ux440ux435ux43bux435}

Источник:
\texttt{Greenhouse/firmware/greenhouse\_mega/greenhouse\_mega.ino}.

Назначение контроллера: прием MQTT-команд, переключение 15 релейных
каналов, публикация подтвержденных состояний и диагностической сводки.

\subsubsection{А.1. Основные
параметры}\label{ux430.1.-ux43eux441ux43dux43eux432ux43dux44bux435-ux43fux430ux440ux430ux43cux435ux442ux440ux44b}

\begin{verbatim}
const char* WIFI_SSID    = "Greenhouse";
const char* WIFI_PASS    = "77777777";

const char* MQTT_HOST    = "192.168.1.112";
const uint16_t MQTT_PORT = 1883;

const char* DEVICE_ID    = "mega-1";
const char* BASE         = "greenhouse";

#define NUM_RELAYS     15
#define SUMMARY_MS     60000UL
#define WIFI_RETRY_MS  5000UL
#define MQTT_RETRY_MS  5000UL
\end{verbatim}

\subsubsection{А.2. Релейные
выходы}\label{ux430.2.-ux440ux435ux43bux435ux439ux43dux44bux435-ux432ux44bux445ux43eux434ux44b}

\begin{verbatim}
uint8_t RELAY_PINS[NUM_RELAYS] = {
  22, 24, 26, 28, 30, 32, 34, 36, 38, 40, 42, 44, 46, 48, 50
};
const bool RELAY_ACTIVE_LOW = false;
\end{verbatim}

\subsubsection{А.3. Публикация состояния
реле}\label{ux430.3.-ux43fux443ux431ux43bux438ux43aux430ux446ux438ux44f-ux441ux43eux441ux442ux43eux44fux43dux438ux44f-ux440ux435ux43bux435}

\begin{verbatim}
void publishRelayState(uint8_t idx) {
  if (!mqtt.connected() || idx >= NUM_RELAYS) return;
  char topic[64];
  snprintf(topic, sizeof(topic), "%s/relay/%u/state", BASE, idx + 1);
  mqtt.publish(topic, relayState[idx] ? "ON" : "OFF", true);
}
\end{verbatim}

\subsubsection{А.4. Обработка
MQTT-команд}\label{ux430.4.-ux43eux431ux440ux430ux431ux43eux442ux43aux430-mqtt-ux43aux43eux43cux430ux43dux434}

Контроллер принимает топики вида
\texttt{greenhouse/relay/\textless{}N\textbar{}all\textgreater{}/\textless{}set\textbar{}get\textgreater{}}.
Команда \texttt{set} меняет состояние, команда \texttt{get} публикует
текущее состояние без изменения реле.

Поддерживаемые команды:

{\def\LTcaptype{none} % do not increment counter
\begin{longtable}[]{@{}ll@{}}
\hline
Команда & Действие \\
\hline
\endhead
\hline
\endlastfoot
\texttt{ON} или \texttt{1} & включить реле \\
\texttt{OFF} или \texttt{0} & выключить реле \\
\texttt{TOGGLE} & переключить состояние \\
\end{longtable}
}

\VkrAppendixSection{Б}{Фрагменты прошивок сенсорных
узлов}\label{ux43fux440ux438ux43bux43eux436ux435ux43dux438ux435-ux431-ux444ux440ux430ux433ux43cux435ux43dux442ux44b-ux43fux440ux43eux448ux438ux432ux43eux43a-ux441ux435ux43dux441ux43eux440ux43dux44bux445-ux443ux437ux43bux43eux432}

\subsubsection{Б.1. Узел температуры и
влажности}\label{ux431.1.-ux443ux437ux435ux43b-ux442ux435ux43cux43fux435ux440ux430ux442ux443ux440ux44b-ux438-ux432ux43bux430ux436ux43dux43eux441ux442ux438}

Источник: \texttt{Greenhouse/firmware/sensor\_th/sensor\_th.ino}.

\begin{verbatim}
const char* DEVICE_ID = "th-1";
#define DHT_PIN   2
#define DHT_TYPE  DHT11
#define READ_INTERVAL_MS   10000UL

snprintf(topicState, sizeof(topicState), "greenhouse/env/%s/state", DEVICE_ID);
snprintf(topicLWT, sizeof(topicLWT), "greenhouse/%s/status", DEVICE_ID);
\end{verbatim}

Формат сообщения:

\begin{verbatim}
{"ts":123456,"device":"th-1","t":24.5,"h":61.0}
\end{verbatim}

\subsubsection{\texorpdfstring{Б.2. Узел
CO\(_2\)}{Б.2. Узел CO\_2}}\label{ux431.2.-ux443ux437ux435ux43b-co_2}

Источник: \texttt{Greenhouse/firmware/sensor\_co2/sensor\_co2.ino}.

\begin{verbatim}
const char* DEVICE_ID = "co2-1";

#define SDA_PIN   4
#define SCL_PIN   5
#define READ_INTERVAL_MS   5000UL

snprintf(topicState, sizeof(topicState), "greenhouse/env/%s/state", DEVICE_ID);
snprintf(topicLWT, sizeof(topicLWT), "greenhouse/%s/status", DEVICE_ID);
\end{verbatim}

Формат сообщения:

\begin{verbatim}
{"ts":123456,"device":"co2-1","co2":820,"t":24.7,"h":60.5}
\end{verbatim}

В прошивке предусмотрена фильтрация некорректных значений CO\(_2\):
значения меньше 0 ppm и больше 10000 ppm не публикуются.

\VkrAppendixSection{В}{MQTT-топики и форматы
сообщений}\label{ux43fux440ux438ux43bux43eux436ux435ux43dux438ux435-ux432-mqtt-ux442ux43eux43fux438ux43aux438-ux438-ux444ux43eux440ux43cux430ux442ux44b-ux441ux43eux43eux431ux449ux435ux43dux438ux439}

{\def\LTcaptype{none} % do not increment counter
\begin{longtable}[]{@{}
  >{\RaggedRight\arraybackslash}p{(\linewidth - 6\tabcolsep) * \real{0.2500}}
  >{\RaggedRight\arraybackslash}p{(\linewidth - 6\tabcolsep) * \real{0.2500}}
  >{\RaggedRight\arraybackslash}p{(\linewidth - 6\tabcolsep) * \real{0.2500}}
  >{\RaggedRight\arraybackslash}p{(\linewidth - 6\tabcolsep) * \real{0.2500}}@{}}
\hline
\begin{minipage}[b]{\linewidth}\RaggedRight
Направление
\end{minipage} & \begin{minipage}[b]{\linewidth}\RaggedRight
Топик
\end{minipage} & \begin{minipage}[b]{\linewidth}\RaggedRight
Payload
\end{minipage} & \begin{minipage}[b]{\linewidth}\RaggedRight
Назначение
\end{minipage} \\
\hline
\endhead
\hline
\endlastfoot
sensor -\textgreater{} broker & \texttt{greenhouse/env/th-1/state} &
JSON \texttt{ts}, \texttt{device}, \texttt{t}, \texttt{h} & температура
и влажность \\
sensor -\textgreater{} broker & \texttt{greenhouse/env/co2-1/state} &
JSON \texttt{ts}, \texttt{device}, \texttt{co2}, \texttt{t}, \texttt{h}
& CO\(_2\), температура, влажность \\
device -\textgreater{} broker &
\texttt{greenhouse/\textless{}device\textgreater{}/status} &
\texttt{online} / \texttt{offline} & доступность узла \\
dashboard -\textgreater{} relay &
\texttt{greenhouse/relay/\textless{}n\textgreater{}/set} & \texttt{ON},
\texttt{OFF}, \texttt{TOGGLE}, \texttt{1}, \texttt{0} & команда на
реле \\
dashboard -\textgreater{} relay & \texttt{greenhouse/relay/all/set} &
\texttt{ON}, \texttt{OFF}, \texttt{TOGGLE}, \texttt{1}, \texttt{0} &
команда на все реле \\
dashboard -\textgreater{} relay &
\texttt{greenhouse/relay/\textless{}n\textgreater{}/get} & любое
сообщение & запрос состояния \\
relay -\textgreater{} broker &
\texttt{greenhouse/relay/\textless{}n\textgreater{}/state} & \texttt{ON}
/ \texttt{OFF} & подтвержденное состояние \\
relay -\textgreater{} broker & \texttt{greenhouse/relay/summary} & JSON
& диагностическая сводка \\
relay -\textgreater{} broker & \texttt{greenhouse/mega-1/status} &
\texttt{online} / \texttt{offline} & состояние контроллера \\
\end{longtable}
}

Пример сводки контроллера:

\begin{verbatim}
{"ts":123456,"device":"mega-1","uptime":3600,"wifi_rc":3,"mqtt_rc":0,"rssi":-61,"relays":{"1":false,"2":false,"3":true}}
\end{verbatim}

Фактические MQTT-логи после запуска оборудования нужно получить по
протоколу \texttt{thesis/test\_protocol.md}.

\VkrAppendixSection{Г}{Структура базы
данных}\label{ux43fux440ux438ux43bux43eux436ux435ux43dux438ux435-ux433-ux441ux442ux440ux443ux43aux442ux443ux440ux430-ux431ux430ux437ux44b-ux434ux430ux43dux43dux44bux445}

Источник: \texttt{Greenhouse/dashboard/app.py}.

\subsubsection{\texorpdfstring{Г.1. Таблица
\texttt{sensor\_log}}{Г.1. Таблица sensor\_log}}\label{ux433.1.-ux442ux430ux431ux43bux438ux446ux430-sensor_log}

{\def\LTcaptype{none} % do not increment counter
\begin{longtable}[]{@{}lll@{}}
\hline
Поле & Тип & Назначение \\
\hline
\endhead
\hline
\endlastfoot
\texttt{id} & INTEGER PRIMARY KEY AUTOINCREMENT & идентификатор
записи \\
\texttt{ts} & TEXT & время записи \\
\texttt{device} & TEXT & идентификатор датчика \\
\texttt{temperature} & REAL & температура \\
\texttt{humidity} & REAL & относительная влажность \\
\texttt{co2} & REAL & концентрация CO\(_2\) \\
\end{longtable}
}

\subsubsection{\texorpdfstring{Г.2. Таблица
\texttt{relay\_log}}{Г.2. Таблица relay\_log}}\label{ux433.2.-ux442ux430ux431ux43bux438ux446ux430-relay_log}

{\def\LTcaptype{none} % do not increment counter
\begin{longtable}[]{@{}lll@{}}
\hline
Поле & Тип & Назначение \\
\hline
\endhead
\hline
\endlastfoot
\texttt{id} & INTEGER PRIMARY KEY AUTOINCREMENT & идентификатор
записи \\
\texttt{ts} & TEXT & время записи \\
\texttt{relay\_id} & INTEGER & номер реле \\
\texttt{state} & INTEGER & состояние 0/1 \\
\end{longtable}
}

\subsubsection{Г.3.
Индексы}\label{ux433.3.-ux438ux43dux434ux435ux43aux441ux44b}

{\def\LTcaptype{none} % do not increment counter
\begin{longtable}[]{@{}lll@{}}
\hline
Индекс & Таблица & Поле \\
\hline
\endhead
\hline
\endlastfoot
\texttt{idx\_sensor\_ts} & \texttt{sensor\_log} & \texttt{ts} \\
\texttt{idx\_relay\_ts} & \texttt{relay\_log} & \texttt{ts} \\
\end{longtable}
}

\subsubsection{Г.4.
SQL-проверка}\label{ux433.4.-sql-ux43fux440ux43eux432ux435ux440ux43aux430}

\begin{verbatim}
select ts, device, temperature, humidity, co2
from sensor_log
order by id desc
limit 5;

select ts, relay_id, state
from relay_log
order by id desc
limit 5;
\end{verbatim}

Фактический вывод этих запросов нужно вставить после запуска
оборудования.

\VkrAppendixSection{Д}{Реле и исполнительные
устройства}\label{ux43fux440ux438ux43bux43eux436ux435ux43dux438ux435-ux434-ux440ux435ux43bux435-ux438-ux438ux441ux43fux43eux43bux43dux438ux442ux435ux43bux44cux43dux44bux435-ux443ux441ux442ux440ux43eux439ux441ux442ux432ux430}

{\def\LTcaptype{none} % do not increment counter
\begin{longtable}[]{@{}rll@{}}
\hline
Реле & Имя по умолчанию & Группа \\
\hline
\endhead
\hline
\endlastfoot
1 & Бактерицидный фильтр & Фильтры \\
2 & 2-й УФ фильтр & Фильтры \\
3 & 1-й УФ фильтр & Фильтры \\
4 & 1-й насос & Насосы \\
5 & 4-й насос & Насосы \\
6 & 2-й насос & Насосы \\
7 & 3-й насос & Насосы \\
8 & 2-й вентилятор & Вентиляторы \\
9 & 4-я заслонка справа & Заслонки \\
10 & 1-й вентилятор & Вентиляторы \\
11 & 3-я заслонка справа & Заслонки \\
12 & 2-я заслонка слева & Заслонки \\
13 & Правый ТЭН & ТЭНы \\
14 & 1-я заслонка слева & Заслонки \\
15 & Левый ТЭН & ТЭНы \\
\end{longtable}
}

\VkrAppendixSection{Е}{Веб-интерфейс
оператора}\label{ux43fux440ux438ux43bux43eux436ux435ux43dux438ux435-ux435-ux432ux435ux431-ux438ux43dux442ux435ux440ux444ux435ux439ux441-ux43eux43fux435ux440ux430ux442ux43eux440ux430}

Необходимо добавить скриншоты после запуска dashboard:

\begin{enumerate}
\def\labelenumi{\arabic{enumi})}
\tightlist
\item
  Главная страница \texttt{/} с состоянием датчиков и реле.
\item
  Страница \texttt{/rules} с примером правила.
\item
  Страница \texttt{/profiles} с примером профиля вентиляции.
\item
  Страница \texttt{/stats} с графиком и экспортом.
\end{enumerate}

Порядок получения скриншотов описан в \texttt{thesis/test\_protocol.md}.

\VkrAppendixSection{Ж}{Схемы}\label{ux43fux440ux438ux43bux43eux436ux435ux43dux438ux435-ux436-ux441ux445ux435ux43cux44b}

Исходники схем находятся в \texttt{thesis/diagrams.md}:

\begin{enumerate}
\def\labelenumi{\arabic{enumi})}
\tightlist
\item
  Общая архитектура системы.
\item
  MQTT-обмен сообщениями.
\item
  Контур правил и профилей.
\item
  Накопление данных и ML-пайплайн.
\item
  SQLite-структура истории.
\end{enumerate}

Перед вставкой в финальную ВКР схемы нужно экспортировать в PNG или SVG.

\VkrAppendixSection{И}{План
ML-эксперимента}\label{ux43fux440ux438ux43bux43eux436ux435ux43dux438ux435-ux438-ux43fux43bux430ux43d-ml-ux44dux43aux441ux43fux435ux440ux438ux43cux435ux43dux442ux430}

Сокращенная постановка:

{\def\LTcaptype{none} % do not increment counter
\begin{longtable}[]{@{}
  >{\RaggedRight\arraybackslash}p{(\linewidth - 2\tabcolsep) * \real{0.5000}}
  >{\RaggedRight\arraybackslash}p{(\linewidth - 2\tabcolsep) * \real{0.5000}}@{}}
\hline
\begin{minipage}[b]{\linewidth}\RaggedRight
Элемент
\end{minipage} & \begin{minipage}[b]{\linewidth}\RaggedRight
Описание
\end{minipage} \\
\hline
\endhead
\hline
\endlastfoot
Исходные данные & \texttt{sensor\_log}, \texttt{relay\_log} \\
Горизонты прогноза & 10, 30 и 60 минут \\
Целевые переменные & температура, влажность, CO\(_2\) \\
Базовые признаки & лаги, скользящие агрегаты, время суток, состояния
реле \\
Baseline & последнее значение, скользящее среднее \\
ML-модели & Random Forest, Gradient Boosting, LSTM/GRU при достаточном
объеме данных \\
Метрики & MAE, RMSE, R2 \\
\end{longtable}
}

Подробный план находится в \texttt{thesis/ml\_experiment\_plan.md}.


\end{document}
